\chapter{Pseudoeffective classes and ideal-valuative realizations}
\label{chap:pseff-mp-dxz}
\index{pseudoeffective class!ideal-valuative realization}
\index{comparison!Mesquita--Piccione--Darvas--Xia--Zhang}

The tropical spectrum $X^{\bethan}$ is attached to the coherent ideals
of $X$, independently of a positivity class.  Positivity enters later and
in two different ways.  Mesquita--Piccione use relatively nef
cohomological models and assume that the background class is K\"ahler.
Darvas--Xia--Zhang instead attach an inverse system of positive-mass test
curves to every pseudoeffective class.  The purpose of this chapter is to
put the inverse-system data on the tropical spectrum without identifying
objects that the existing theories do not identify.

Throughout, $X$ is a connected compact K\"ahler manifold of dimension
$n$, $\theta$ is a smooth closed real $(1,1)$-form, and
\[
                    \alpha:=\{\theta\}
       \in H^{1,1}_{\mathrm{BC}}(X,\mathbf R)
\tag{71.1}\label{eq:pseff-background}
\]
is pseudoeffective.  The smoothness of $X$ is used for multiplier ideals,
their valuative characterization, and strong openness.  The DXZ space
$\PSH^{\NA}(X,\theta)$ is the inverse system constructed in
\Cref{chap:transcendental-metrics}.  The MP space $X^{\bethan}$, its
semivaluations, and its divisorial locus were constructed in
\Cref{chap:tropical-spectra-complex,chap:mp-dual-complexes}.

There are three conclusions to keep separate.
\begin{enumerate}
\item The DXZ divisorial analytification is canonical and order-reflecting;
      this is the content of \cite[Definition~3.11 and Lemma~4.2]
      {Xia23Operations}.
\item Evaluating the same multiplier ideals at an arbitrary semivaluation
      gives a canonical raw function on $X^{\bethan}$.  This extension is
      formal.  Its upper semicontinuity is proved below from finite ideal
      data, but that regularity alone supplies no positivity.
\item For every K\"ahler class on a compact K\"ahler manifold, the raw
      transform is an order isomorphism onto the full MP cone, with no
      projectivity, rationality, or finite-energy hypothesis. The proof
      first identifies the finite-energy loci, then uses dominated
      divisorial approximation and explicit truncations. In the smooth
      projective ample N\'eron--Severi setting, the MP and
      Boucksom--Jonsson cones also agree. The remaining boundaries concern
      pseudoeffective classes outside these proved regimes.
\end{enumerate}
The limitations in the third clause are part of the hypotheses of the
comparison theorem, not merely notational restrictions.

\section{The obstruction carried by the generic fibre}
\label{sec:pseff-generic-fibre}

An unchanged extension of the MP model definition would ask for a
cohomological test configuration
$\pi:\cX\to\mathbf P^1$ and a relatively nef class
$\mathcal B\in H^{1,1}_{\mathrm{BC}}(\cX,\mathbf R)$ whose restriction to
each noncentral fibre is $\alpha$.  The following elementary obstruction
explains why this cannot treat a non-nef class.

\begin{proposition}\label{prop:pseff-relative-nef-obstruction}
Let $\pi:\cX\to\mathbf P^1$ be a cohomological test configuration in
the sense of \cite[Definitions~2.5.1 and 4.1.2]{MP24}; thus $\cX$ is a
normal compact K\"ahler space, but need not be smooth.  Let
$\mathcal B\in H^{1,1}_{\mathrm{BC}}(\cX,\mathbf R)$ be relatively nef.
Assume that, under the canonical identification
$\cX_t\simeq X$ for $t\ne0$,
\[
                        \mathcal B|_{\cX_t}=\alpha.
\tag{71.2}\label{eq:pseff-generic-restriction}
\]
Then $\alpha$ is nef.  More generally, any vertically corrected class,
on $\cX$ or on a modification that is an isomorphism away from the
central fibre, which is relatively nef and has generic-fibre restriction
$\alpha$, still forces $\alpha$ to be nef.
\end{proposition}

\begin{proof}
By the definition used in \cite[Definition~4.1.2]{MP24}, relative
nefness implies that the restriction of $\mathcal B$ to each fibre is
nef.  Since the equivariant trivialization identifies $\cX_t$ with the
smooth manifold $X$ for $t\ne0$,
\eqref{eq:pseff-generic-restriction} therefore says directly that
$\alpha$ is nef.  Equivalently, if $\mathcal H$ is a relatively K\"ahler
class and one uses the approximation definition of relative nefness, then
restriction to a noncentral fibre gives the K\"ahler classes
\[
        \alpha+\varepsilon\,\mathcal H|_{\cX_t}
              \in H^{1,1}_{\mathrm{BC}}(X,\mathbf R)
\tag{71.3}\label{eq:pseff-nef-restriction}
\]
for every $\varepsilon>0$, and letting $\varepsilon\downarrow0$ gives
the same conclusion.

If $D$ is a vertical divisor, then
$[D]|_{\cX_t}=0$ for $t\ne0$, so a corrected class
$\mathcal B+[D]$ has the same generic-fibre restriction.  If this
corrected class is relatively nef, the first paragraph applied to
$\mathcal B+[D]$ proves the assertion.  Similarly, if
$\mu:\cX'\to\cX$ is an isomorphism away from $\cX_0$, then
$\mu:\cX'_t\to\cX_t$ is an isomorphism for $t\ne0$.  Every exceptional
vertical class restricts to zero on $\cX'_t$; hence the restriction of
$\mu^*\mathcal B$, with any such correction, is still $\alpha$.
If the corrected pullback class on $\cX'$ is relatively nef, applying
the first paragraph to that corrected relatively nef class on $\cX'$
proves the last assertion.  No claim is made that an arbitrary vertical
correction preserves relative nefness.  No smoothness of either total
space is used.
\end{proof}

This is a direct consequence of the MP definition, rather than a new
positivity theorem.  The proof works in the full normal MP
test-configuration category.  The
generic fibre is already one of the compact analytic subspaces tested by
relative nefness; a curve-only substitute would not be sufficient on an
arbitrary nonprojective K\"ahler manifold.

\begin{example}[A nef and big boundary class]\label{ex:pseff-nef-big}
Let $p:\widetilde{\mathbf P}^2\to\mathbf P^2$ be the blowup at a point,
let $H=p^*c_1(\mathcal O_{\mathbf P^2}(1))$, and let $E$ be the
exceptional curve.  Then
\[
                         H^2=1,\qquad H\cdot E=0,
                         \qquad E^2=-1.
\tag{71.4}\label{eq:pseff-blowup-intersections}
\]
Thus $H$ is nef and big, but it is not K\"ahler because it has zero
intersection with $E$.  The product test configuration with class
$\operatorname{pr}_1^*H$ is relatively nef, so
\Cref{prop:pseff-relative-nef-obstruction} permits this model.  It does not
assert that the rest of the MP K\"ahler theory automatically extends to
$H$.
\end{example}

\begin{example}[A big non-nef class]\label{ex:pseff-big-nonnef}
On the same surface, fix $0<a<1$ and put
$\alpha_a=H+aE$.  Since $H$ is big and $aE$ is effective,
$\alpha_a$ is big.  On the other hand,
\[
                         \alpha_a\cdot E=-a<0,
\tag{71.5}\label{eq:pseff-big-nonnef-intersection}
\]
so $\alpha_a$ is not nef.  Its surface Zariski decomposition is
$\alpha_a=H+aE$, with positive part $H$ and negative part $aE$;
in particular $\vol(\alpha_a)=H^2=1$.
\Cref{prop:pseff-relative-nef-obstruction} shows that no relatively nef
cohomological test configuration can restrict to $\alpha_a$ on its
noncentral fibres.  The DXZ inverse system nevertheless exists, since
$\alpha_a$ is pseudoeffective.
\end{example}

\begin{remark}\label{rem:pseff-obstruction-scope}
The proposition rules out the unchanged relatively nef MP model signature
over a non-nef class.  It does not rule out a different future construction
using movable products, positive parts, or lower-dimensional data.
\end{remark}

\section{Ideal-valuative numbers and the two equivalence relations}
\label{sec:pseff-ideal-values}

Let $v\in X^{\bethan}$.  For a coherent ideal sheaf $\mathfrak a$,
we use the convention
\[
 v(\mathfrak a):=\min\{v(f):f\text{ is a local generator of }
 \mathfrak a\text{ at the centre of }v\}.
\tag{71.6}\label{eq:pseff-ideal-value}
\]
The value is independent of the generators.  It lies in
$[0,+\infty]$, and it is allowed to be infinite.  If $\varphi$ is a
quasi-psh function on $X$, set
\[
 \nu_v^{\mathcal I}(\varphi)
 :=\lim_{m\to\infty}\frac{1}{m}
      v\bigl(\mathcal I(m\varphi)\bigr).
\tag{71.7}\label{eq:pseff-ideal-valuative-number}
\]

\begin{proposition}\label{prop:pseff-ideal-valuative-basic}
The limit in \eqref{eq:pseff-ideal-valuative-number} exists in
$[0,+\infty]$, and
\[
 \nu_v^{\mathcal I}(\varphi)
 =\sup_{m\geq1}\frac{1}{m}
      v\bigl(\mathcal I(m\varphi)\bigr).
\tag{71.8}\label{eq:pseff-fekete-formula}
\]
For $s>0$, one has
$\nu_{sv}^{\mathcal I}(\varphi)=s\nu_v^{\mathcal I}(\varphi)$.  If
$v=r\ord_E$, where $E$ is a prime divisor over $X$ and
$r>0$, then
\[
             \nu_{r\ord_E}^{\mathcal I}(\varphi)
             =r\nu_E(\varphi).
\tag{71.9}\label{eq:pseff-divisorial-lelong}
\]
\end{proposition}

\begin{proof}
The subadditivity theorem for multiplier ideals gives
\[
 \mathcal I((m+\ell)\varphi)
 \subseteq\mathcal I(m\varphi)\mathcal I(\ell\varphi).
\tag{71.10}\label{eq:pseff-multiplier-subadditivity}
\]
Since inclusion reverses ideal values and $v(\mathfrak a\mathfrak b)
=v(\mathfrak a)+v(\mathfrak b)$, the sequence
$b_m:=v(\mathcal I(m\varphi))$ is superadditive.  Fekete's lemma,
valid also with the value $+\infty$, proves
\eqref{eq:pseff-fekete-formula}.  For $s>0$, homogeneity follows first at
every finite stage from $(sv)(\mathfrak a)=s v(\mathfrak a)$, and then by
taking the supremum; no product $0\cdot(+\infty)$ is formed.  Finally, the
asymptotic divisorial order of the
multiplier ideals is the generic Lelong number:
\[
 \lim_{m\to\infty}m^{-1}\ord_E\mathcal I(m\varphi)
 =\nu_E(\varphi).
\]
This is \cite[Lemma~B.4]{BBJ21}, and multiplying by $r$ proves
\eqref{eq:pseff-divisorial-lelong}.
\end{proof}

The next statement records precisely how much singularity data
\eqref{eq:pseff-ideal-valuative-number} remembers.

\begin{proposition}\label{prop:pseff-i-equivalence-valuative}
For quasi-psh functions $\varphi$ and $\psi$ on the smooth manifold
$X$, the following conditions are equivalent:
\begin{enumerate}
\item $\mathcal I(m\varphi)=\mathcal I(m\psi)$ for every integer
      $m\geq1$;
\item $\nu_E(\varphi)=\nu_E(\psi)$ for every prime divisor $E$
      over $X$;
\item $\nu_v^{\mathcal I}(\varphi)=\nu_v^{\mathcal I}(\psi)$ for
      every divisorial semivaluation $v=r\ord_E$.
\end{enumerate}
Consequently these conditions are equivalent to
$\mathcal I$-equivalence in the sense of
\cite[Definition~2.7]{Xia23Operations}.
\end{proposition}

\begin{proof}
The first condition implies the third by
\eqref{eq:pseff-ideal-valuative-number}, and the third is equivalent to
the second by \eqref{eq:pseff-divisorial-lelong}.  Suppose now that the
second condition holds.  The valuative characterization of multiplier
ideals in \cite[Theorem~A]{BFJ08} characterizes the upper-regularized
ideal $\mathcal I_+(m\varphi)$ by inequalities involving precisely the
numbers $\nu_E(\varphi)$ and the log discrepancies $A_X(E)$.
Hence
$\mathcal I_+(m\varphi)=\mathcal I_+(m\psi)$ for every $m$.
Strong openness \cite{GZ15} gives
$\mathcal I_+(m\varphi)=\mathcal I(m\varphi)$, and likewise for
$\psi$, proving the first condition.  Finally,
\cite[Definition~2.3]{Xia23Operations} defines the
$\mathcal I$-envelope by equality of all the integer multiplier ideals;
its normalized fixed points are therefore classified by the same data.
This is exactly the equivalence relation of
\cite[Definition~2.7]{Xia23Operations}.
\end{proof}

\begin{warning}{The $P$-versus-$\mathcal I$ firewall}
\label{warn:pseff-p-i-firewall}
The transform below records multiplier ideals, hence
$\mathcal I$-equivalence.  It must not be used to identify arbitrary
$P$-equivalence classes.  The two relations agree at an
$\mathcal I$-good positive-mass potential, where by definition
\[
                  P_\eta[\varphi]
                  =P_\eta[\varphi]_{\mathcal I},
\tag{71.11}\label{eq:pseff-i-good-firewall}
\]
but no such identity is asserted outside that locus.  In particular, a
future theory sensitive to movable mass may have to retain more data than
the function in \eqref{eq:pseff-ideal-valuative-number}.
\end{warning}

\section{The transform of a positive-mass test curve}
\label{sec:pseff-one-curve}

Let $\eta$ be a smooth closed real $(1,1)$-form with big class, and
let
$\Gamma=(\Gamma_\tau)\in\PSH^{\NA}(X,\eta)_{>0}$ be a DXZ
$\mathcal I$-model positive-mass test curve.  Write
$c:=\Gamma_{\max}$.  Thus the interior levels
$\Gamma_\tau$, $\tau<c$, are normalized
$\mathcal I$-model potentials.  Define
\[
 \mathfrak T_\Gamma(v)
 :=\sup_{\tau<c}
       \{\tau-\nu_v^{\mathcal I}(\Gamma_\tau)\}
 \in[-\infty,c]
 \qquad (v\in X^{\bethan}).
\tag{71.12}\label{eq:pseff-test-curve-transform}
\]
For $b>0$, it is convenient to homogenize this function by
\[
 \widehat{\mathfrak T}_\Gamma(v,b)
 :=b\mathfrak T_\Gamma(b^{-1}v)
 =\sup_{\tau<c}
       \{b\tau-\nu_v^{\mathcal I}(\Gamma_\tau)\}.
\tag{71.13}\label{eq:pseff-homogeneous-transform}
\]

\begin{proposition}\label{prop:pseff-one-curve-transform}
The transform in \eqref{eq:pseff-test-curve-transform} has the following
properties.
\begin{enumerate}
\item \(\mathfrak T_\Gamma(0)=c\), and
      \(\mathfrak T_\Gamma\leq c\) on \(X^{\bethan}\).
\item For every \(v\), the restriction of
      \(s\mapsto\mathfrak T_\Gamma(sv)\) to \((0,+\infty)\) is either
      identically \(-\infty\), or finite, convex, and nonincreasing.
\item For \(\lambda>0\),
      \[
       \widehat{\mathfrak T}_\Gamma(\lambda v,\lambda b)
       =\lambda\widehat{\mathfrak T}_\Gamma(v,b).
      \tag{71.14}\label{eq:pseff-joint-homogeneity}
      \]
\item If \(C\in\mathbf R\), then
      \(\mathfrak T_{\Gamma+C}=\mathfrak T_\Gamma+C\).
\item If \(\Gamma\leq\Gamma'\) level by level, then
      \(\mathfrak T_\Gamma\leq\mathfrak T_{\Gamma'}\).
\end{enumerate}
No upper semicontinuity or MP plurisubharmonicity on all of
\(X^{\bethan}\) is included in the statement.
\end{proposition}

\begin{proof}
Ideal values are nonnegative, so each term in
\eqref{eq:pseff-test-curve-transform} is at most \(\tau<c\).  This proves
\(\mathfrak T_\Gamma\leq c\).  At the zero semivaluation all ideal values
vanish, whence
\(\mathfrak T_\Gamma(0)=\sup_{\tau<c}\tau=c\).

By homogeneity in \Cref{prop:pseff-ideal-valuative-basic},
\[
 \mathfrak T_\Gamma(sv)
 =\sup_{\tau<c}
      \{\tau-s\nu_v^{\mathcal I}(\Gamma_\tau)\}.
\tag{71.15}\label{eq:pseff-raywise-affine-supremum}
\]
If every coefficient that occurs is \(+\infty\), the right side is
identically \(-\infty\) for \(s>0\).  Otherwise the terms with infinite
coefficient can be discarded, and the right side is a finite-valued
supremum of affine functions of \(s\).  It is therefore convex; since all
slopes are
nonpositive, it is nonincreasing.  The second expression
in \eqref{eq:pseff-homogeneous-transform} and the same homogeneity give
\eqref{eq:pseff-joint-homogeneity} immediately.

By \cite[Definition~4.3]{Xia23Operations},
\((\Gamma+C)_\tau=\Gamma_{\tau-C}\).  Substituting
\(\rho=\tau-C\) in \eqref{eq:pseff-test-curve-transform} proves the fourth
claim.  Finally, \(\Gamma_\tau\leq\Gamma'_\tau\) makes
\(\Gamma_\tau\) at least as singular as \(\Gamma'_\tau\).  Monotonicity of multiplier ideals therefore gives
\(\nu_v^{\mathcal I}(\Gamma_\tau)\geq
\nu_v^{\mathcal I}(\Gamma'_\tau)\).  Taking the supremum in
\eqref{eq:pseff-test-curve-transform} proves the fifth claim.
\end{proof}

\begin{example}[The constant ray]
Assume that \(\eta\) itself is a K\"ahler form.  Fix \(a\in\mathbf R\) and
put \(\Gamma_\tau=0\) for \(\tau\leq a\), with
\(\Gamma_\tau=-\infty\) for \(\tau>a\).  Its inverse Legendre transform is
the ray \(u_t=at\).  Since
\(\mathcal I(m\Gamma_\tau)=\mathcal O_X\) for \(\tau\leq a\),
\[
                         \mathfrak T_\Gamma(v)=a
                         \quad\text{for every }v\in X^{\bethan}.
\tag{71.16}\label{eq:pseff-constant-transform}
\]
Thus endpoint translation becomes addition of a constant, exactly as in
\Cref{prop:pseff-one-curve-transform}.
\end{example}

\begin{example}[A fixed logarithmic pole]
\label{ex:pseff-fixed-log-pole}
Let $X=\mathbf P^1$, let $\eta$ be the Fubini--Study form normalized
by $\int_X\eta=1$, and let $s_p$ vanish simply at $p$.  Fix
$0<\lambda<1$, and let
\[
 \varphi:=P_\eta[\lambda\log|s_p|^2]_{\mathcal I}.
\tag{71.17}\label{eq:pseff-log-pole-potential}
\]
This is an $\mathcal I$-model potential of nonpluripolar mass
$1-\lambda>0$, and it has generic Lelong number $\lambda$ at $p$.
The relative test curve that equals $\varphi$ for $\tau\leq0$ and
equals $-\infty$ for $\tau>0$ has
\[
 \mathfrak T_\Gamma(r\ord_p)=-r\lambda,
 \qquad r>0.
\tag{71.18}\label{eq:pseff-log-pole-transform}
\]
Indeed, \eqref{eq:pseff-divisorial-lelong} makes every term in the
supremum equal to $\tau-r\lambda$, whose supremum for $\tau<0$ is
$-r\lambda$.  If a semivaluation is supported at $p$, then the
multiplier ideals eventually restrict to zero on its support and the same
transform can take the value $-\infty$.  This is why the full target of
\eqref{eq:pseff-test-curve-transform} is extended real-valued.
More concretely, let $v_p$ be evaluation at $p$, so a germ nonzero at
$p$ has value zero and a germ vanishing there has value $+\infty$.
Then $q_N=N\ord_p\preceq v_p$ for every $N$, and $q_N\to v_p$
on every coherent ideal of $\mathbf P^1$. The transforms satisfy
$\mathfrak T_\Gamma(q_N)=-N\lambda\to-\infty
 =\mathfrak T_\Gamma(v_p)$. This is a concrete dominated divisorial
approximation of a point with proper support.

This curve is bounded relative to its initial model potential $\varphi$.
Its initial mass is $m=1-\lambda$, and every interior level has that
same mass. Hence the normalized relative energy is zero, whereas the
class-normalized DXZ integral is $-\infty$:
\[
 \frac1{1-\lambda}\int_{-\infty}^{0}
       ((1-\lambda)-(1-\lambda))\,\rd\tau=0,
 \qquad
 \int_{-\infty}^{0}((1-\lambda)-1)\,\rd\tau=-\infty.
 \label{eq:pseff-relative-energy-counterexample}
\]
Thus finite relative energy need not place a curve in the finite-energy
space used by \Cref{thm:pseff-kahler-finite-energy-link}.
\end{example}

\section{The radial convention and rational Gauss weights}
\label{sec:pseff-radial}

This section records a source-level formula from the projective part of
\cite{DXZ23}.  Accordingly, until the next section we assume in addition
that $X$ is projective.  No compact-K\"ahler extension of the formula is
used later.  We fix the radial variable by
\[
                         t=-\log|z|^2.
\tag{71.19}\label{eq:pseff-radial-convention}
\]
The inverse Legendre transform of $\Gamma$ is
\[
 u_t:=\uscsup[\tau<c]\{\Gamma_\tau+t\tau\},
 \qquad t\geq0.
\tag{71.20}\label{eq:pseff-inverse-legendre}
\]
On $X\times\Delta^*$, write
$U(x,z):=u_{-\log|z|^2}(x)$.  If $q>c$, then
\[
 W_q(x,z):=U(x,z)+q\log|z|^2
          =u_t(x)-qt
\tag{71.21}\label{eq:pseff-shifted-radial-potential}
\]
extends as a quasi-psh function across $X\times\{0\}$.

Let $v=r\ord_E$, with $r\in\mathbf Q_{>0}$, and let
$b\in\mathbf Q_{>0}$.  Its weight-$b$ Gauss extension is
\[
 \sigma_{v,b}\!\left(\sum_{j\geq0}f_jz^j\right)
 :=\min_j\{v(f_j)+bj\}.
\tag{71.22}\label{eq:pseff-weighted-gauss}
\]
It is a rational divisorial valuation over $X\times\Delta$, after a
common resolution and a ramified base change.  In particular the generic
Lelong number $\nu_{\sigma_{v,b}}(W_q)$ is defined.  Our normalization is
\[
                  \nu_{\sigma_{v,b}}(\log|z|^2)=b.
\tag{71.23}\label{eq:pseff-gauss-log-normalization}
\]

\begin{proposition}\label{prop:pseff-weighted-radial-formula}
With the hypotheses and conventions above,
\[
 \widehat{\mathfrak T}_\Gamma(v,b)
 =qb-\nu_{\sigma_{v,b}}(W_q).
\tag{71.24}\label{eq:pseff-weighted-radial-formula}
\]
The right side is independent of the chosen $q>c$.
The statement is restricted to rational divisorial $v$ and rational
$b>0$; it asserts no formula for an arbitrary nondivisorial
semivaluation.
\end{proposition}

\begin{proof}
Translate the test curve by $-q$.  By
\Cref{prop:pseff-one-curve-transform}, its endpoint is $c-q<0$, its
dual ray is $u_t-qt$, and its complexification is precisely $W_q$.
The divisorial Gauss identity \cite[Proposition~3.1]{DXZ23}, applied to
the divisorial valuation
$b^{-1}v$, gives
\[
 \mathfrak T_\Gamma(b^{-1}v)-q
 =-\nu_{\sigma_{b^{-1}v,1}}(W_q).
\tag{71.25}\label{eq:pseff-unit-gauss-step}
\]
The definitions give an exact equality of valuations
\[
                 \sigma_{v,b}=b\sigma_{b^{-1}v,1}.
\tag{71.26}\label{eq:pseff-gauss-scaling}
\]
Generic Lelong numbers are homogeneous under positive rescaling of a
divisorial valuation.  Multiplying \eqref{eq:pseff-unit-gauss-step} by
$b$, and using \eqref{eq:pseff-homogeneous-transform} and
\eqref{eq:pseff-gauss-scaling}, proves
\eqref{eq:pseff-weighted-radial-formula}.

If $q'>c$, then
$W_{q'}=W_q+(q'-q)\log|z|^2$.  By
\eqref{eq:pseff-gauss-log-normalization},
\[
 \nu_{\sigma_{v,b}}(W_{q'})
 =\nu_{\sigma_{v,b}}(W_q)+(q'-q)b.
\]
Thus $q'b-\nu_{\sigma_{v,b}}(W_{q'})$ equals the right side of
\eqref{eq:pseff-weighted-radial-formula}, proving independence of $q$.
\end{proof}

\section{Passage to a pseudoeffective class}
\label{sec:pseff-inverse-system}

Let now
$\boldsymbol\Gamma=(\Gamma^{\theta+\omega})_\omega\in
\PSH^{\NA}(X,\theta)$ be a compatible DXZ inverse-system metric as in
\Cref{def:tna-pshna}.  The common endpoint is denoted by
$c=\boldsymbol\Gamma_{\max}$.  Choose any K\"ahler perturbation
$\omega$, and set
\[
 \mathfrak T_{\boldsymbol\Gamma}(v)
 :=\sup_{\tau<c}
   \{\tau-\nu_v^{\mathcal I}
       (\Gamma^{\theta+\omega}_\tau)\}.
\tag{71.27}\label{eq:pseff-inverse-system-transform}
\]

\begin{proposition}\label{thm:pseff-perturbation-independent-transform}
The right side of \eqref{eq:pseff-inverse-system-transform} is independent
of \(\omega\).  It therefore defines a canonical map
\[
 \mathfrak T_\theta:
 \PSH^{\NA}(X,\theta)
 \longrightarrow[-\infty,+\infty)^{X^{\bethan}},
 \qquad
 \boldsymbol\Gamma\longmapsto
 \mathfrak T_{\boldsymbol\Gamma}.
\tag{71.28}\label{eq:pseff-canonical-transform-map}
\]
For every \(\boldsymbol\Gamma\), its image satisfies all five conclusions
of \Cref{prop:pseff-one-curve-transform}.  In particular,
\(\mathfrak T_{\boldsymbol\Gamma}(0)=\boldsymbol\Gamma_{\max}\).
The proposition makes no assertion that the raw function in
\eqref{eq:pseff-inverse-system-transform} is upper semicontinuous or MP
psh at nondivisorial points.
\end{proposition}

\begin{proof}
Let \(\omega\) and \(\omega'\) be two perturbations.  Choose a third
perturbation dominating both in the directed system.  It is enough to
compare either original component with its image under one transition.
By \eqref{eq:tna-transition}, at every \(\tau<c\) that image is
\[
 P_{\theta+\omega'}
 [\Gamma^{\theta+\omega}_\tau]_{\mathcal I}
\tag{71.29}\label{eq:pseff-transition-level}
\]
with the forms relabelled according to the order convention of
\eqref{eq:tna-kahler-order}.  The \(\mathcal I\)-envelope is
\(\mathcal I\)-equivalent to its input.  By
\Cref{prop:pseff-i-equivalence-valuative}, this gives the exact equality
\[
 \mathcal I\!\left(mP_{\theta+\omega'}
 [\Gamma^{\theta+\omega}_\tau]_{\mathcal I}\right)
 =\mathcal I(m\Gamma^{\theta+\omega}_\tau),
 \qquad m\in\mathbf Z_{>0}.
\tag{71.30}\label{eq:pseff-transition-ideal-equality}
\]
Taking ideal values, dividing by \(m\), and passing to the limit proves
that \(\nu_v^{\mathcal I}\) is unchanged for every
\(v\in X^{\bethan}\).  The endpoint is common by
\Cref{prop:tna-common-terminal}.  Thus
\eqref{eq:pseff-inverse-system-transform} is independent of the
component.  All remaining assertions follow by applying
\Cref{prop:pseff-one-curve-transform} to any one component.
\end{proof}

The proposition is the arbitrary-semivaluation rewriting of the transition
formula in \cite[Definition~3.3]{DXZ23}.  Its divisorial restriction is
exactly the analytification of
\cite[Definition~3.11]{Xia23Operations}.  The only extra observation is
that equality of coherent multiplier ideals may be evaluated at any
semivaluation; no new positivity assertion is involved.

The fact that the target of \eqref{eq:pseff-canonical-transform-map} is a
space of functions, rather than an MP psh cone, is deliberate.  On the
divisorial locus no new inversion theorem is needed: the exact statement is
already in the literature.

\begin{theorem}\label{thm:pseff-divisorial-inversion}
For \(v=r\ord_E\), with \(r>0\), one has
\[
 \mathfrak T_{\boldsymbol\Gamma}(v)
 =\sup_{\tau<c}\{\tau-r\nu_E(
       \Gamma^{\theta+\omega}_\tau)\}
 =\boldsymbol\Gamma^{\mathrm{an}}(v),
\tag{71.31}\label{eq:pseff-divisorial-analytification}
\]
where the last term is the DXZ analytification of
\cite[Definition~3.11]{Xia23Operations}.  Consequently the divisorial
restriction of \(\mathfrak T_\theta\), together with its value at the
zero semivaluation, is injective and reflects order:
\[
 \boldsymbol\Gamma\leq\boldsymbol\Gamma'
 \quad\Longleftrightarrow\quad
 \mathfrak T_{\boldsymbol\Gamma}(v)
 \leq\mathfrak T_{\boldsymbol\Gamma'}(v)
 \quad\text{for }v=0\text{ and every divisorial }v,
\tag{71.32}\label{eq:pseff-order-reflection}
\]
where the value at \(v=0\) records the endpoint inequality.
\end{theorem}

\begin{proof}
Equation \eqref{eq:pseff-divisorial-lelong} changes the defining formula
\eqref{eq:pseff-inverse-system-transform} into the first expression in
\eqref{eq:pseff-divisorial-analytification}.  This is precisely
\cite[formula~(3.7)]{Xia23Operations}.  The equivalence
\eqref{eq:pseff-order-reflection} is
\cite[Lemma~4.2]{Xia23Operations}; its proof recovers the levelwise order
from all divisorial Legendre values and the maximality of the normalized
\(\mathcal I\)-model levels.  Applying the equivalence in both directions
shows that equality of the divisorial restrictions forces equality of the
DXZ metrics.  Finally,
\(\mathfrak T_{\boldsymbol\Gamma}(0)=c\) records the terminal parameter,
including when \(\dim X=0\).
\end{proof}

\begin{remark}\label{rem:pseff-injection-not-bijection}
\Cref{thm:pseff-divisorial-inversion} is an imported divisorial
determination statement, not a
bijection onto a previously defined MP cone.  An intrinsic
characterization of the image in
$[ -\infty,+\infty)^{X^{\bethan}}$, especially at nondivisorial points,
is one of the open boundaries recorded at the end of the chapter.
\end{remark}

\subsection{Upper semicontinuity from finite ideal data}
\label{sec:pseff-raw-usc}

The next argument resolves the upper-semicontinuity question for the raw
ideal-valuative transform.  Its finite approximants are continuous functions,
but no positivity is attached to them.

\begin{lemma}\label{lem:p142-monotone-parameter-supremum}
Let $Y$ be a topological space, let $c\in\mathbf R$, and, for each
$\tau<c$, let $a_\tau:Y\to[0,+\infty]$ be lower semicontinuous.
Assume that $a_s\leq a_t$ whenever $s\leq t<c$. Then
\[
 F(y):=\sup_{\tau<c}\{\tau-a_\tau(y)\}
 \in[-\infty,c]
\]
is upper semicontinuous. For every dense subset $D\subset(-\infty,c)$,
\[
 F(y)=\sup_{\tau\in D}\{\tau-a_\tau(y)\}
 \quad(y\in Y).
 \tag{71.33a}
 \label{eq:p142-dense-parameter-formula}
\]
No continuity in $\tau$, countability assumption on $Y$, or finiteness
assumption on $a_\tau(y)$ is required.
\end{lemma}

\begin{theorem}\label{thm:p142-raw-transform-usc}
Let $X$ be a complex manifold, and give the tropical spectrum
$X^{\bethan}$ of its coherent ideal sheaves the topology of pointwise
convergence on ideals. Let $c\in\mathbf R$ and let
$(\Gamma_\tau)_{\tau<c}$ be a
decreasing family of quasi-plurisubharmonic functions on $X$, none
identically $-\infty$ on a connected component. Put
\[
 \nu_v^{\mathcal I}(\Gamma_\tau)
 :=\sup_{m\geq1}\frac1m v\bigl(\mathcal I(m\Gamma_\tau)\bigr),
 \qquad
 F_\Gamma(v):=\sup_{\tau<c}
       \{\tau-\nu_v^{\mathcal I}(\Gamma_\tau)\}.
 \tag{71.33b}
 \label{eq:p142-raw-transform}
\]
The supremum defining $\nu_v^{\mathcal I}$ equals the limit as
$m\to\infty$. Then $F_\Gamma$ is upper semicontinuous on all of
$X^{\bethan}$, including semivaluations with nonzero support. Formula
\eqref{eq:p142-dense-parameter-formula} holds for $F_\Gamma$ at every
semivaluation.

More precisely, for $N,M\in\mathbf Z_{>0}$ set
\[
 \tau_{N,j}:=c-N+\frac jN \qquad(0\leq j<N^2).
 \tag{71.33c}
 \label{eq:p142-mesh}
\]
\[
 A_{\tau,M}(v):=\max_{1\leq m\leq M}
       \frac1m v\bigl(\mathcal I(m\Gamma_\tau)\bigr).
 \tag{71.33d}
 \label{eq:p142-finite-ideal-values}
\]
\[
 B_{N,M}(v):=\frac1N+
   \max\left\{c-N,
       \max_{0\leq j<N^2}
        \bigl(\tau_{N,j}-A_{\tau_{N,j},M}(v)\bigr)\right\}.
 \tag{71.33e}
 \label{eq:p142-upper-approximants}
\]
\[
 f_k(v):=\min_{1\leq N,M\leq k}B_{N,M}(v).
 \tag{71.33f}
 \label{eq:p142-decreasing-sequence}
\]
Every $B_{N,M}$ and $f_k$ is a bounded, real-valued continuous function
given by a finite expression in negative ideal evaluations $-v(\mathfrak a)$,
real translations, positive rational scaling, minima and maxima. Moreover,
\[
 f_k\searrow F_\Gamma\quad\hbox{pointwise on }X^{\bethan}.
 \tag{71.33g}
 \label{eq:p142-exact-sequential-limit}
\]
The approximants are not asserted to be MP plurisubharmonic or
semipositive.
\end{theorem}

\begin{proof}[Proof of \cref{lem:p142-monotone-parameter-supremum}]
For $N\geq1$ put $a=c-N$, and partition $[a,c)$ into the $N^2$
half-open intervals
\[
 [\tau_{N,j},\tau_{N,j}+1/N),\qquad
 \tau_{N,j}=c-N+j/N,\quad 0\leq j<N^2.
\]
Only left endpoints are sampled; the family is never evaluated at $c$.
Set
\[
 F_N:=\max\{c-N,F\},\qquad
 G_N:=\max\left\{c-N,
      \max_{0\leq j<N^2}
       (\tau_{N,j}-a_{\tau_{N,j}})\right\}.
\]
Every sampled term occurs in the supremum defining $F$, so $G_N\leq F_N$.
If $\tau<c-N$, nonnegativity gives $\tau-a_\tau\leq c-N$.
Otherwise there is $j$ with
$\tau_{N,j}\leq\tau<\tau_{N,j}+1/N$. Monotonicity in the parameter
gives, in extended-real order,
\[
 \tau-a_\tau
 \leq \tau-a_{\tau_{N,j}}
 \leq \tau_{N,j}-a_{\tau_{N,j}}+1/N
 \leq G_N+1/N.
\]
This also holds when one or both of the ideal-type quantities are infinite:
only a finite real number minus $+\infty$ occurs. Taking the supremum and
including the floor proves
\[
             G_N\leq F_N\leq G_N+1/N.
 \tag{71.33h}
 \label{eq:p142-grid-sandwich}
\]

Each $G_N$ is a finite-valued upper-semicontinuous function, since
$-a_{\tau_{N,j}}$ is upper semicontinuous. The sandwich gives the exact
identity
\[
                  F=\inf_{N\geq1}(G_N+1/N).
 \tag{71.33i}
 \label{eq:p142-usc-infimum}
\]
Indeed, $F\leq F_N\leq G_N+1/N$, and
$G_N+1/N\leq F_N+1/N$ tends pointwise to $F$ as $N\to\infty$.
If $F(y)=-\infty$, then $F_N(y)=c-N$, so the same conclusion holds
at that point. An arbitrary infimum of upper-semicontinuous functions is
upper semicontinuous: its strict sublevel set is the union of the strict
sublevel sets of its terms. This proves the first assertion on an
arbitrary topological space.

For the dense-parameter assertion, fix $y$ and $\tau<c$ with
$a_\tau(y)<+\infty$. Given $\varepsilon>0$, choose
$s\in D\cap(\tau-\varepsilon,\tau)$. Then
\[
 s-a_s(y)\geq s-a_\tau(y)
              >\tau-a_\tau(y)-\varepsilon.
\]
Thus the supremum over $D$ is at least every finite term in the full
supremum, after letting $\varepsilon$ tend to zero. Terms equal to
$-\infty$ impose no further condition. If all terms are $-\infty$,
both suprema are $-\infty$. The opposite inequality follows from
$D\subset(-\infty,c)$, proving
\eqref{eq:p142-dense-parameter-formula}.
\end{proof}

\begin{proof}[Proof of \cref{thm:p142-raw-transform-usc}]
Multiplier ideal sheaves of quasi-psh functions on a complex manifold
are coherent.  To apply the psh subadditivity theorem, work on a relatively
compact coordinate neighbourhood and write
\(\Gamma_\tau=\varphi_\tau+\rho_\tau\), where
\(\varphi_\tau\) is psh and \(\rho_\tau\) is smooth.  For every
\(r>0\), the factor \(\exp(-r\rho_\tau)\) is bounded above and below
by positive constants on a smaller relatively compact neighbourhood.
Consequently
\[
             \mathcal I(r\Gamma_\tau)
             =\mathcal I(r\varphi_\tau)
\]
there.  Applying \cite[Theorem~2.6(ii)]{DEL00} to
\(m\varphi_\tau\) and \(\ell\varphi_\tau\), and then using the
same bounded-perturbation identity, gives
\[
 \mathcal I((m+\ell)\Gamma_\tau)
 \subseteq\mathcal I(m\Gamma_\tau)
               \mathcal I(\ell\Gamma_\tau).
\]
Ideal values reverse inclusion and add under multiplication. Hence
$b_m=v(\mathcal I(m\Gamma_\tau))$ is a nonnegative extended-real
superadditive sequence. If some $b_q=+\infty$, superadditivity implies
$b_m=+\infty$ for every $m\geq q$: use $m=q+r$ for $r\geq1$.
Otherwise, for a fixed $q$, writing $m=kq+r$, with $0\leq r<q$, gives
$b_m\geq kb_q$ (put $b_0=0$). Therefore
\[
 \liminf_{m\to\infty}\frac{b_m}{m}\geq\frac{b_q}{q}.
\]
Taking the supremum over $q$ and comparing with the immediate upper
bound proves the limit assertion, including an infinite supremum. This
is the multiplier-ideal argument recorded in
\cref{prop:pseff-ideal-valuative-basic}.

By the topology of the tropical spectrum, the function
$v\mapsto v(\mathfrak a)$ is continuous into $[0,+\infty]$ with its
order topology for every coherent ideal $\mathfrak a$; see
\cite[Definition~1.3.1]{MP24} and \eqref{eq:ts-evaluation}.
It follows that
\[
 v\longmapsto\nu_v^{\mathcal I}(\Gamma_\tau)
   =\sup_{m\geq1}\frac1m
          v(\mathcal I(m\Gamma_\tau))
\]
is lower semicontinuous. If $s\leq t<c$, then
$\Gamma_s\geq\Gamma_t$. Local integrability immediately implies
\[
 \mathcal I(m\Gamma_t)\subseteq\mathcal I(m\Gamma_s),
 \qquad
 \nu_v^{\mathcal I}(\Gamma_s)
       \leq\nu_v^{\mathcal I}(\Gamma_t).
\]
For the first inclusion, multiplication of the weight by the bounded
factor associated to an additive local constant would also suffice;
the given pointwise order already yields it directly. No assertion
about continuity of multiplier ideals in the parameter is used.
The hypotheses of \cref{lem:p142-monotone-parameter-supremum} are now
verified, and that lemma proves upper semicontinuity and the dense-set
formula at every semivaluation.

We prove the stronger explicit approximation assertion separately.
For fixed $N,M$, all ideal evaluations in
\eqref{eq:p142-upper-approximants} are continuous extended-real
functions. Their displayed finite maxima and positive scalings are
continuous in the order topology; subtracting from the finite
$\tau_{N,j}$ gives a continuous function into $[-\infty,\tau_{N,j}]$.
The outer maximum with $c-N$ then yields an ordinary real-valued
continuous function. In particular,
\[
        c-N+1/N\leq B_{N,M}\leq c+1/N.
\]
This establishes boundedness without invoking compactness of $X^{\bethan}$.

As $M$ increases, $A_{\tau,M}\nearrow\nu^{\mathcal I}(\Gamma_\tau)$.
For the finite number of sampled parameters, taking the maximum commutes
with these decreasing limits after changing the sign. To see this even
at $-\infty$, if every limiting term is below a chosen real threshold,
then every term is eventually below it; a maximum of finitely many
indices works for all terms. The fixed finite floor is retained throughout.
Consequently, with $G_N$ as in the preceding proof,
\[
           B_{N,M}\searrow G_N+1/N\quad(M\to\infty).
\]
The sandwich \eqref{eq:p142-grid-sandwich} gives
\[
 F_\Gamma\leq B_{N,M},\qquad
 \inf_{N,M\geq1}B_{N,M}
      =\inf_{N\geq1}(G_N+1/N)=F_\Gamma.
 \tag{71.33j}
 \label{eq:p142-double-infimum}
\]
Every $f_k$ is a finite minimum of bounded continuous real-valued
functions. The squares $\{1,\ldots,k\}^2$ increase and exhaust
$\mathbf Z_{>0}^2$, so their minima decrease to the double infimum
in \eqref{eq:p142-double-infimum}. This is
\eqref{eq:p142-exact-sequential-limit}.

The construction is explicit. It does not extract a sequence from an
arbitrary net on a nonmetrizable space. Nor does a minimum of these
continuous functions acquire a positivity property merely because it
is formed from ideal values. No positive approximation is asserted.
\end{proof}

\begin{corollary}\label{cor:p142-pseudoeffective-raw-usc}
Let $X$ be a connected compact K\"ahler manifold, let $\theta$ be a
smooth closed real $(1,1)$-form with pseudoeffective class, and let
$\boldsymbol\Gamma\in\PSH^{\mathrm{NA}}(X,\theta)$ be a DXZ
inverse-system metric. Its perturbation-independent raw ideal-valuative
transform is upper semicontinuous on $X^{\bethan}$ and has the decreasing
continuous finite-ideal approximation in
\eqref{eq:p142-decreasing-sequence}--\eqref{eq:p142-exact-sequential-limit},
formed from any one perturbation component. Every dense parameter set
computes the same transform at all semivaluations.

This assertion requires neither projectivity nor rationality, bigness,
nefness, or positive volume of the unperturbed class. It asserts no MP
positivity, positive approximation, Monge--Amp\`ere construction, or
identification of $P$-equivalence with $\mathcal I$-equivalence.
\end{corollary}

\begin{proof}[Proof of \cref{cor:p142-pseudoeffective-raw-usc}]
Fix one K\"ahler perturbation component
$\Gamma=(\Gamma^{\theta+\omega}_\tau)_{\tau<c}$, where
$c=\boldsymbol\Gamma_{\max}$. The DXZ test-curve definition supplies
a decreasing family of quasi-psh functions with finite terminal
parameter, and its interior levels are not identically $-\infty$;
see \cite[Section~2.2, Definition~3.3]{DXZ23}.
\Cref{thm:p142-raw-transform-usc} therefore applies to this component.

The identification with the raw inverse-system transform uses only the
already established transition identity, not an MP comparison theorem.
Under a transition the interior level is replaced by its
$\mathcal I$-model envelope in the larger class. This preserves every
integer multiplier ideal by \cite[Proposition~2.18(ii)]{DX22}; the
standing hypotheses of its Section~2.4 allow an arbitrary compact
K\"ahler manifold and a pseudoeffective background class. A common
dominating component compares any
two choices, and their terminal parameters coincide. Evaluating those
same ideals at any semivaluation proves that both the raw transform and
the explicit expressions \eqref{eq:p142-finite-ideal-values}--
\eqref{eq:p142-decreasing-sequence} are unchanged. These are precisely
the ideal equalities in
\cref{thm:pseff-perturbation-independent-transform}; the source
transition is \cite[formula~(28), Definition~3.3]{DXZ23}.

All arguments take place on a perturbed component. No potential of
positive mass in the unperturbed class is chosen, so zero-volume
pseudoeffective classes are included. The proof does not construct
relatively nef cohomological models, a Monge--Amp\`ere operator, or an
energy. Those are distinct questions.
\end{proof}

\begin{example}\label{ex:p142-infinite-ideal-value}
On $X=\mathbf P^1$, choose a point $p$, a smooth Hermitian metric on
$\mathcal O(p)$, and its canonical section $s_p$. Fix $\lambda>0$ and
put $\psi=\lambda\log|s_p|^2$. The family $\Gamma_\tau=\psi$ for
$\tau<c$ satisfies the elementary hypotheses of
\cref{thm:p142-raw-transform-usc}; no assertion that this particular
representative is an $\mathcal I$-model potential is needed.

Writing $\mathfrak m_p$ for the ideal sheaf of $p$, one has
\[
       \mathcal I(m\psi)=\mathfrak m_p^{\lfloor m\lambda\rfloor}.
\]
Indeed, near $p$ a coordinate $z$ expresses $\psi$ as
$\lambda\log|z|^2$ plus a smooth function. A germ of order $q$ is
integrable with the weight $\mathrm e^{-m\psi}$ exactly when
$q-m\lambda>-1$, which is equivalent to $q\geq\lfloor m\lambda\rfloor$.
Away from $p$ the ideal is the unit ideal. It follows that
\[
 \nu_v^{\mathcal I}(\psi)=\lambda v(\mathfrak m_p),
 \qquad F_\Gamma(v)=c-\lambda v(\mathfrak m_p).
\]
At a semivaluation with $v(\mathfrak m_p)=+\infty$, each sufficiently
large multiplier-ideal value is infinite, so the displayed formulas give
$F_\Gamma(v)=-\infty$. For the smaller integers with
$\lfloor m\lambda\rfloor=0$, the ideal is the unit ideal and its value
is zero; no expression $0\cdot(+\infty)$ is used. Thus the infinite-value
case in the theorem records actual support information and cannot be
discarded by restricting to divisorial valuations.
\end{example}

\section{Divisorial piecewise-linear approximation}
\label{sec:pseff-ideal-approximation}

The transform can be approached by finite ideal data, but the order of the
two limits matters.  Fix a perturbation component and a finite set
$T\subset(-\infty,c)$, assumed nonempty. For $M\geq1$, put
\[
 a_{\tau,M}(v)
 :=\max_{1\leq m\leq M}
      \frac1m v\bigl(\mathcal I(m\Gamma_\tau)\bigr),
\qquad
 \mathfrak T_{T,M}(v)
 :=\max_{\tau\in T}\{\tau-a_{\tau,M}(v)\}.
\tag{71.36}\label{eq:pseff-finite-ideal-approximant}
\]

\begin{proposition}\label{prop:pseff-divisorial-pl-approximation}
Each $\mathfrak T_{T,M}$ is an extended-real-valued finite ideal
expression on $X^{\bethan}$.  On the zero-support valuation locus---in
particular, on every dual complex and at every rational divisorial
point---it is a real piecewise-linear function.  For fixed $T$, it
decreases as $M$ increases, and
\[
 \lim_{M\to\infty}\mathfrak T_{T,M}(v)
 =\max_{\tau\in T}
    \{\tau-\nu_v^{\mathcal I}(\Gamma_\tau)\}.
\tag{71.37}\label{eq:pseff-finite-grid-limit}
\]
Let $D$ be a countable dense subset of $( -\infty,c)$, and let
$T_1\subset T_2\subset\cdots$ be finite subsets whose union is $D$.
Then, at every semivaluation $v\in X^{\bethan}$,
\[
 \mathfrak T_{\boldsymbol\Gamma}(v)
 =\lim_{N\to\infty}\lim_{M\to\infty}
       \mathfrak T_{T_N,M}(v).
\tag{71.38}\label{eq:pseff-iterated-pl-limit}
\]
The iterated limit is an extended-real identity on the whole spectrum.
Neither the finite approximants nor this limit are declared MP psh by this
approximation argument.
\end{proposition}

\begin{proof}
For a coherent ideal $\mathfrak a$, the function
$v\mapsto-v(\mathfrak a)$ is piecewise linear on every dual complex
carrying $\mathfrak a$, and these functions are compatible under
refinement; see \Cref{chap:tropical-spectra-complex}.  Thus
$-m^{-1}v(\mathcal I(m\Gamma_\tau))$ is real piecewise linear there.
Finite minima and maxima of piecewise-linear functions are piecewise
linear after passing to a common subdivision.  Since
\[
 \tau-a_{\tau,M}(v)
 =\min_{1\leq m\leq M}
   \left\{\tau-\frac1m
             v\bigl(\mathcal I(m\Gamma_\tau)\bigr)\right\},
\]
the second formula in \eqref{eq:pseff-finite-ideal-approximant} is
piecewise linear on that locus.  At a semivaluation with nonzero support,
an appearing ideal may have value $+\infty$; the same finite formula then
takes values in $[ -\infty,+\infty)$, which is why no globally
real-valued PL assertion is made.

By \eqref{eq:pseff-fekete-formula},
$a_{\tau,M}(v)$ increases to
$\nu_v^{\mathcal I}(\Gamma_\tau)$.  For a finite set $T$, the maximum
of finitely many decreasing extended-real sequences decreases to the
maximum of their limits.  This proves
\eqref{eq:pseff-finite-grid-limit}.

At an arbitrary semivaluation $v$, the quantities
$\nu_v^{\mathcal I}(\Gamma_\tau)$ are nonnegative and nondecreasing in
$\tau$. By \Cref{lem:p142-monotone-parameter-supremum}, the dense set
$D$ computes their defining transform, even when every term is
$-\infty$. The maxima over the increasing finite sets $T_N$ increase
to this dense-set supremum. Combining that fact with
\eqref{eq:pseff-finite-grid-limit} proves
\eqref{eq:pseff-iterated-pl-limit} at every semivaluation.
\end{proof}

\begin{example}[The two limits for a logarithmic pole]
\label{ex:pseff-finite-ideal-grid}
Use \Cref{ex:pseff-fixed-log-pole} with $\lambda=2/5$, and evaluate at
$v=(3/2)\ord_p$. For each interior parameter,
$\mathcal I(m\Gamma_\tau)=\mathfrak m_p^{\lfloor2m/5\rfloor}$.
With $T=\{-1,-1/2,-1/4\}$ and $M=4$, the finite slope is
\[
 a_{\tau,4}(v)=\frac32
     \max\left\{0,0,\frac13,\frac14\right\}=\frac12,
 \qquad \mathfrak T_{T,4}(v)=-\frac34.
\]
Keeping $T$ fixed and letting $M\to\infty$ gives
$\max T-(3/2)(2/5)=-17/20$. Then grids approaching the endpoint zero
increase this value to the exact transform $-3/5$.
Thus the first limit decreases at a fixed grid, while the second enlarges
that grid. Neither step supplies a positive model approximation.
\end{example}

\begin{remark}\label{rem:pseff-pl-limit-firewall}
The finite-stage expressions in \eqref{eq:pseff-finite-ideal-approximant}
are a
computational device.  To prove that their limit is an MP psh function
one would need a decreasing net of \emph{relatively nef} model
functions.  Piecewise linearity alone supplies no such positivity, and
interchanging the grid and multiplier-ideal limits is not justified.
\end{remark}

\section{The proved comparison domain and the projective boundary}
\label{sec:pseff-comparison}

We first state the comparison that is available without projectivity but
with a finite-energy hypothesis.  Let $\alpha$ be K\"ahler, let
$\omega\in\alpha$ be a K\"ahler form, and normalize
$V_\omega=0$. In the finite-energy correspondence below, DXZ finite
energy uses the class volume $V_\alpha=\int_X\omega^n$, as in
\cite[formula~(43), Section~3.4]{DXZ23}; the initial level is $0$.
This differs from finite relative energy based at a singular model
potential of smaller mass. The radial energy here is divided by
$V_\alpha$, matching the probability normalization of the MP source.
Write
$\mathcal E^{1,\NA}_{\mathrm{DXZ}}(X,\omega)$ for the DXZ
finite-energy test curves and $\mathcal E^1_{\mathrm{MP}}(\alpha)$ for
the finite-energy cone of \Cref{chap:mp-psh}.

There is an essential convention change between the sources.
Darvas--Xia--Zhang use $t=-\log|z|^2$ and normalize a generic Lelong
number against $\log|z|^2$, while \cite{MP24} uses
$s=-\log|z|$ and the ordinary-norm convention.  Passing from an
ordinary-norm weight to the book's squared-norm potential doubles both the
potential and the logarithmic denominator.  Hence the numerical
non-Archimedean value is unchanged; in particular no factor of two is to be
inserted in \eqref{eq:pseff-finite-energy-divisorial-link}.

\begin{theorem}\label{thm:pseff-kahler-finite-energy-link}
There is a bijection
\[
 \mathcal E^{1,\NA}_{\mathrm{DXZ}}(X,\omega)
 \xrightarrow{\ \sim\ }
 \mathcal E^1_{\mathrm{MP}}(\alpha).
\tag{71.39}\label{eq:pseff-finite-energy-bijection}
\]
It sends a test curve \(\Gamma\) to the MP potential \(\phi\) determined
on rational divisorial points by
\[
 \phi(p^{-1}\ord_E)
 =\mathfrak T_\Gamma(p^{-1}\ord_E),
 \qquad E\text{ prime},\quad p\in\mathbf Q_{>0}.
\tag{71.40}\label{eq:pseff-finite-energy-divisorial-link}
\]
Equivalently, both sides encode the same maximal radial asymptotic data.
For any chosen \(u_0\in\mathcal E^1(X,\omega)\),
\Cref{thm:mp-maximal-ray} selects the unique maximal representative with
initial value \(u_0\).  The radial and non-Archimedean energies agree under
\eqref{eq:pseff-finite-energy-bijection}.
\end{theorem}

\begin{proof}
By \cite[Theorem~3.17]{DXZ23}, inverse Legendre transform is a bijection
from finite-energy \(\mathcal I\)-maximal test curves to the corresponding
finite-energy rays emanating from \(V_\omega=0\), and it identifies radial
energy with DXZ non-Archimedean energy.  Applying the radial construction
of \cite[Theorem~5.1.4 and Section~5.2]{MP24} gives an element of
\(\mathcal E^1_{\mathrm{MP}}(\alpha)\).  Conversely,
\cite[Theorem~5.1.7]{MP24}, restated as
\Cref{thm:mp-maximal-ray}, applied with initial value \(0\), supplies the
maximal ray whose Ross--Witt Nystr\"om transform recovers the DXZ test
curve.  This proves \eqref{eq:pseff-finite-energy-bijection}.  A different
initial value changes the maximal representative, not its radial
non-Archimedean data.

For \(v=p^{-1}\ord_E\), \cite[Theorem~5.1.4 and Section~5.2]{MP24}
defines the MP value as the negative generic Lelong number of the
complexified ray and identifies it with the associated DXZ divisorial
potential.  Equation \eqref{eq:pseff-divisorial-analytification} identifies
that potential with \(\mathfrak T_\Gamma(v)\).  The convention comparison
preceding the theorem shows that the two source normalizations give the
same coefficient.  This proves
\eqref{eq:pseff-finite-energy-divisorial-link}.  Uniqueness of the MP psh
extension follows from \cite[Theorem~4.3.8]{MP24}.  Finally, equality of
energies is the last assertion of \cite[Theorem~3.17]{DXZ23}, together
with the MP slope formula \cite[Theorem~5.3.4]{MP24}.
\end{proof}

\begin{remark}\label{rem:pseff-finite-energy-raw-function}
The theorem characterizes the MP potential by the divisorial restriction
of the raw ideal-valuative transform.  It does not assert, on a general
compact K\"ahler manifold, that the two functions have already been proved
equal at every nondivisorial point.  The next subsection supplies that
equality in the finite-energy K\"ahler domain. The subsequent truncation
theorem extends it to the full K\"ahler cones.
\end{remark}

\subsection{Domination-preserving divisorial approximation}
\label{sec:pseff-dominated-density}

Ordinary density of rational divisorial valuations does not control the
ideal order needed by the raw transform.  The following refinement does.
Write $q\preceq v$ when
\[
 q(\mathfrak a)\leq v(\mathfrak a)
 \qquad\text{for every coherent ideal }\mathfrak a\subseteq\mathcal O_X.
\tag{71.41a}\label{eq:p143-ideal-order}
\]
Here $X^{\divis}$ follows the convention of
\cite[Definition~2.3.1]{MP24}: it contains the trivial semivaluation $0$,
as well as the nonnegative rational multiples of divisor orders on proper
bimeromorphic models.  Thus the zero-weight case below is part of the same
divisorial locus, rather than an exception to the statement.

\begin{theorem}\label{thm:p143-dominated-divisorial-density}
For every \(v\in X^{\bethan}\),
\[
 v\in\overline{\{q\in X^{\divis}:q\preceq v\}}.
\tag{71.41b}\label{eq:p143-dominated-divisorial-density}
\]
Equivalently, there is a net \((q_\lambda)\) of rational divisorial
points with \(q_\lambda\preceq v\) and \(q_\lambda\to v\) in the
pointwise topology on coherent ideals.  This includes semivaluations with
proper support and with value \(+\infty\) on a nonzero ideal.

Consequently, let \(\alpha\) be a K\"ahler class, let
\(F:X^{\bethan}\to[-\infty,c]\) be upper semicontinuous and antitone for
\eqref{eq:p143-ideal-order}, and let \(\phi\in\PSH_{\mathrm{MP}}(\alpha)\).
If \(F=\phi\) on \(X^{\divis}\), then \(F=\phi\) on all of
\(X^{\bethan}\).
\end{theorem}

\begin{proof}
Let \(\pi:Y\to X\) be a proper bimeromorphic morphism of connected
compact K\"ahler manifolds.  Pullback induces a continuous map
\(\pi^{\bethan}:Y^{\bethan}\to X^{\bethan}\).  Its image is closed,
because \(Y^{\bethan}\) is compact and \(X^{\bethan}\) is Hausdorff.  It
contains \(X^{\divis}\) by birational invariance of divisorial valuations
\cite[Remark~2.3.2]{MP24}; hence it equals \(X^{\bethan}\) by
\Cref{thm:mp-divisorial-density}.  In particular every \(v\) has a lift
\(\widetilde v\in Y^{\bethan}\).

We first record a domination construction on a simple-normal-crossing
model.  Let \(B_1,\ldots,B_k\) be distinct smooth SNC divisors on \(Y\)
with \(\widetilde v(\mathcal O_Y(-B_i))>0\), and choose finite rational
numbers
\[
 0<a_i\leq\widetilde v(\mathcal O_Y(-B_i)).
\tag{71.41c}\label{eq:p143-dominated-weights}
\]
Let \(Z\) be the central variety of \(\widetilde v\).  By
\cite[Lemma~2.1.1]{MP24}, its ideal \(\mathcal I_Z\) is maximal among
coherent ideals of positive value, so \(Z\subseteq B_i\) for every
\(i\).  Let \(W\) be the connected component of
\(B_1\cap\cdots\cap B_k\) containing \(Z\).  The monomial valuation
\(q\) along \(W\) with weights \((a_i)\) is independent of the point of
\(W\) by \cite[Proposition~2.3.4]{MP24}, and belongs to \(Y^{\divis}\)
by \cite[Remark~2.3.5]{MP24}.

We claim that \(q\preceq\widetilde v\).  For \(s>0\), set
\[
 \mathfrak c_s=
 \sum_{\substack{\beta\in\mathbf Z_{\geq0}^k\\
                   \sum_i a_i\beta_i\geq s}}
 \mathcal O_Y\left(-\sum_i\beta_iB_i\right).
\tag{71.41d}\label{eq:p143-monomial-ideal}
\]
After deleting redundant generators, this is a finite coherent sum: its
exponent set is upward closed, and each minimal exponent has
\(\beta_i\leq\lceil s/a_i\rceil\).  In extended-real order,
\[
 \widetilde v(\mathfrak c_s)
 =\min_{\beta}\sum_{i:\,\beta_i>0}
       \beta_i\widetilde v(\mathcal O_Y(-B_i))
 \geq\min_{\beta}\sum_i\beta_i a_i\geq s.
\tag{71.41e}\label{eq:p143-monomial-value-bound}
\]
The restriction to positive exponents avoids any expression
\(0\cdot(+\infty)\).

Let \(\mathfrak j\) be a nonzero coherent ideal and take
\(0<s<q(\mathfrak j)\).  In SNC coordinates transverse to \(W\), the
condition \(q(\mathfrak j)>s\) says exactly that every transverse Taylor
coefficient of weighted degree less than \(s\) vanishes identically along
\(W\).  There are only finitely many such exponents, and therefore
\(\mathfrak j\subseteq\mathfrak c_s\) on a neighbourhood of \(W\).
Put \(Q=(\mathfrak j+\mathfrak c_s)/\mathfrak c_s\) and
\(\mathfrak k=\operatorname{Ann}(Q)\).  Then \(\mathfrak k\) is the
unit ideal near \(W\) and \(\mathfrak k\mathfrak j\subseteq
\mathfrak c_s\) globally.  It has \(\widetilde v(\mathfrak k)=0\): a
positive value would force \(\mathfrak k\subseteq\mathcal I_Z\), which
is impossible near \(Z\subseteq W\).  Thus
\[
 \widetilde v(\mathfrak j)
 =\widetilde v(\mathfrak k\mathfrak j)
 \geq\widetilde v(\mathfrak c_s)\geq s.
\]
Letting \(s\uparrow q(\mathfrak j)\) proves the claim; the cases
\(q(\mathfrak j)=0\) and \(\mathfrak j=0\) are immediate.

Now fix finitely many nonzero coherent ideals
\(\mathfrak a_1,\ldots,\mathfrak a_\ell\) on \(X\).  Simultaneous
analytic principalization yields a proper modification \(\pi:Y\to X\),
with \(Y\) smooth compact K\"ahler, such that
\[
 \mathfrak a_j\mathcal O_Y=
 \mathcal O_Y(-D_j),\qquad
 D_j=\sum_{i=1}^r d_{ji}B_i,\qquad d_{ji}\in\mathbf Z_{\geq0},
\tag{71.41f}\label{eq:p143-simultaneous-principalization}
\]
where the distinct \(B_i\) form an SNC divisor.  Lift \(v\), and put
\(w_i=\widetilde v(\mathcal O_Y(-B_i))\).  Retain exactly the indices
with \(w_i>0\), choose rational \(a_i\) as in
\eqref{eq:p143-dominated-weights}, and give the other components weight
zero.  If no index is retained, take \(q_Y=0\); otherwise the preceding
construction gives a rational divisorial \(q_Y\preceq\widetilde v\).  In
both cases,
\(q_X:=\pi^{\bethan}(q_Y)\in X^{\divis}\) and \(q_X\preceq v\).
For every \(j\),
\[
 v(\mathfrak a_j)=\sum_{i:\,d_{ji}>0}d_{ji}w_i,\qquad
 q_X(\mathfrak a_j)=\sum_{i:\,d_{ji}>0}d_{ji}a_i.
\tag{71.41g}\label{eq:p143-principalized-values}
\]
Indeed, a zero-weight component cannot contain \(W\): otherwise
\(\mathcal O_Y(-B_i)\subseteq\mathcal I_Z\), forcing \(w_i>0\).
Its monomial value is consequently zero, as used in the second formula.

For every finite positive \(w_i\), let \(a_i\) increase rationally to
\(w_i\); for \(w_i=+\infty\), let \(a_i\) tend to infinity.  The finite
sums in \eqref{eq:p143-principalized-values} then converge in
\([0,+\infty]\).  More explicitly, direct the pairs \((S,N)\), where
\(S\) is a finite set of nonzero coherent ideals and \(N\geq1\), by
inclusion in the first coordinate and the usual order in the second.  For
each pair, principalize the ideals in \(S\) and choose the displayed
weights so that every finite value on \(S\) is approximated from below
within \(1/N\), while every infinite value on \(S\) exceeds \(N\).
The resulting rational divisorial point \(q_{S,N}\preceq v\) belongs to
each corresponding basic pointwise neighbourhood.  Hence
\((q_{S,N})\) converges to \(v\), with no order relation required between
two separately chosen divisorial points.

For the final statement, choose the dominated net.  Antitonicity and
divisorial agreement give
\(F(v)\leq F(q_\lambda)=\phi(q_\lambda)\).  Upper semicontinuity of
\(\phi\) gives \(F(v)\leq\phi(v)\), also when \(\phi(v)=-\infty\), by
testing each finite upper threshold.  Conversely,
\cite[Theorem~4.3.8]{MP24} applied to the upper-semicontinuous function
\(F\) promotes the divisorial inequality \(\phi\leq F\) to all of
\(X^{\bethan}\).  Hence \(F=\phi\).
\end{proof}

\begin{corollary}\label{cor:p143-all-point-kahler-comparison}
Let \(\alpha\) be a K\"ahler class and let
\(\Gamma\in\mathcal E^{1,\NA}_{\mathrm{DXZ}}(X,\omega)\).  If \(\phi\)
is its corresponding MP finite-energy potential under
\eqref{eq:pseff-finite-energy-bijection}, then
\[
 \mathfrak T_\Gamma(v)=\phi(v)
 \qquad(v\in X^{\bethan}).
\tag{71.41h}\label{eq:p143-all-point-kahler-comparison}
\]
\end{corollary}

\begin{proof}
By \Cref{thm:p142-raw-transform-usc}, \(\mathfrak T_\Gamma\) is upper
semicontinuous.  If \(q\preceq v\), every multiplier-ideal evaluation at
\(q\) is at most its evaluation at \(v\).  The defining formula
\eqref{eq:pseff-test-curve-transform} therefore gives
\(\mathfrak T_\Gamma(q)\geq\mathfrak T_\Gamma(v)\).  The divisorial
identity \eqref{eq:pseff-finite-energy-divisorial-link} covers the
\emph{nonzero} rational multiples of divisor orders; the convention for
\(X^{\divis}\) fixed before
\Cref{thm:p143-dominated-divisorial-density} also includes the trivial
semivaluation \(0\), and there the identity holds too:
\(\mathfrak T_\Gamma(0)=\boldsymbol\Gamma_{\max}\) by
\Cref{prop:pseff-one-curve-transform}\,(1), while
\(\phi(0)=\sup_{X^{\bethan}}\phi=\boldsymbol\Gamma_{\max}\) by the maximum principle
at the trivial semivaluation together with \eqref{eq:pseff-finite-energy-bijection}.
Hence \(\mathfrak T_\Gamma=\phi\) on all of \(X^{\divis}\), and
\Cref{thm:p143-dominated-divisorial-density} applies, proving
\eqref{eq:p143-all-point-kahler-comparison}.
\end{proof}

\begin{corollary}\label{cor:p141-finite-energy-raw-image}
For a compact K\"ahler manifold and a K\"ahler class \(\alpha\), the
raw ideal-valuative transform induces a bijection
\[
 \mathcal E^{1,\NA}_{\mathrm{DXZ}}(X,\omega)
 \xrightarrow{\ \mathfrak T\ }
 \mathcal E^1_{\mathrm{MP}}(\alpha).
\tag{71.41i}\label{eq:p141-finite-energy-raw-image}
\]
Thus, in the finite-energy K\"ahler domain, its image is intrinsically
the finite-energy MP cone, viewed as functions on \(X^{\bethan}\).
\end{corollary}

\begin{proof}
For \(\Gamma\in\mathcal E^{1,\NA}_{\mathrm{DXZ}}(X,\omega)\), let
\(\phi\in\mathcal E^1_{\mathrm{MP}}(\alpha)\) be its image under
\Cref{thm:pseff-kahler-finite-energy-link}.  Then
\Cref{cor:p143-all-point-kahler-comparison} gives
\(\mathfrak T_\Gamma=\phi\) on all of \(X^{\bethan}\).  Conversely,
every \(\phi\in\mathcal E^1_{\mathrm{MP}}(\alpha)\) has a unique
preimage \(\Gamma\) under that theorem, and the same comparison identifies
its raw transform with \(\phi\).  The two maps are therefore inverse
bijective maps.
\end{proof}

\begin{theorem}\label{thm:p143-mp-section}
Let \(X\) be a connected compact K\"ahler manifold, let
\(\alpha=\{\omega\}\) be a K\"ahler class.  Then the raw
ideal-valuative transform has a unique section
\[
 \mathcal S_\alpha:\PSH_{\mathrm{MP}}(\alpha)
        \longrightarrow\PSH^{\NA}_{\mathrm{DXZ}}(X,\omega),
 \qquad
 \mathfrak T_{\mathcal S_\alpha(\phi)}=\phi
 \quad\text{on }X^{\bethan}.
\]
The section preserves and reflects pointwise order.

More precisely, let \((\phi_i)_{i\in I}\) be any decreasing net in
\(\PL(X^{\bethan})\cap\PSH_{\mathrm{MP}}(\alpha)\) with pointwise limit
\(\phi\not\equiv-\infty\).  For every \(i\), let
\(\boldsymbol\Gamma^i\in
\mathcal E^{1,\NA}_{\mathrm{DXZ}}(X,\omega)\) be the object corresponding
to \(\phi_i\) under \eqref{eq:pseff-finite-energy-bijection}.  Then the
\(\boldsymbol\Gamma^i\) form a decreasing net,
\[
                 \inf_i\boldsymbol\Gamma^i_{\max}>-\infty,
\]
and
\[
                 \mathcal S_\alpha(\phi)
                 =\inf_i\boldsymbol\Gamma^i.
\]
The infimum and its raw transform are independent of the chosen model net.
Thus every MP \(\alpha\)-psh function, including one not known to have
finite energy, belongs to the image of the raw transform.
\end{theorem}

\begin{proof}
Fix \(\phi\in\PSH_{\mathrm{MP}}(\alpha)\).  By
\Cref{def:mp-general-psh}, there is a decreasing net
\[
 \phi_i\in\PL(X^{\bethan})\cap\PSH_{\mathrm{MP}}(\alpha),
 \qquad \phi_i\downarrow\phi.
\]
Every rational PL MP-psh function has finite model energy, so
\eqref{eq:pseff-finite-energy-bijection} supplies a unique
\(\boldsymbol\Gamma^i\in
\mathcal E^{1,\NA}_{\mathrm{DXZ}}(X,\omega)\).  By
\Cref{cor:p143-all-point-kahler-comparison},
\[
                 \mathfrak T_{\boldsymbol\Gamma^i}=\phi_i
                 \quad\text{on }X^{\bethan}.
\]
If \(i\geq j\), then \(\phi_i\leq\phi_j\) at the positive rational
divisorial points and at the trivial semivaluation.  The order-reflection
statement \eqref{eq:pseff-order-reflection} therefore gives
\(\boldsymbol\Gamma^i\leq\boldsymbol\Gamma^j\).  Thus the DXZ net is
decreasing; no order-compatibility hypothesis has been added.

The trivial semivaluation belongs to \(X^{\divis}\), and
\Cref{thm:mp-psh-divisorial} gives \(\phi(0)>-\infty\).  The endpoint
identity in \Cref{thm:pseff-perturbation-independent-transform} yields
\[
 \boldsymbol\Gamma^i_{\max}
   =\mathfrak T_{\boldsymbol\Gamma^i}(0)
   =\phi_i(0)\downarrow\phi(0)>-\infty.
\]
Hence the decreasing infimum
\(\boldsymbol\Gamma:=\inf_i\boldsymbol\Gamma^i\) exists in
\(\PSH^{\NA}_{\mathrm{DXZ}}(X,\omega)\).  The standard endpoint extension
in \Cref{def:operations-infimum} gives
\[
 \boldsymbol\Gamma_{\max}
   =\inf_i\boldsymbol\Gamma^i_{\max}
   =\phi(0).
\]
Therefore \Cref{thm:pseff-perturbation-independent-transform} gives
\(\mathfrak T_{\boldsymbol\Gamma}(0)=\phi(0)\).  At every positive
rational divisorial point \(q\), \Cref{prop:pseff-operation-formulas} gives
\[
 \mathfrak T_{\boldsymbol\Gamma}(q)
   =\inf_i\mathfrak T_{\boldsymbol\Gamma^i}(q)
   =\inf_i\phi_i(q)
   =\phi(q).
\]
Thus the two functions agree on all of \(X^{\divis}\).

By
\Cref{thm:p142-raw-transform-usc,cor:p142-pseudoeffective-raw-usc},
\(\mathfrak T_{\boldsymbol\Gamma}\) is upper semicontinuous.  It is
antitone for the ideal order: if \(q\preceq v\), then every multiplier-ideal
evaluation at \(q\) is at most the corresponding evaluation at \(v\), and
the defining supremum for the raw transform reverses this inequality.
The comparison clause of \Cref{thm:p143-dominated-divisorial-density},
applied to \(F=\mathfrak T_{\boldsymbol\Gamma}\) and the MP-psh function
\(\phi\), now promotes the divisorial equality to
\[
                 \mathfrak T_{\boldsymbol\Gamma}=\phi
                 \quad\text{on }X^{\bethan}.
\]

Set \(\mathcal S_\alpha(\phi)=\boldsymbol\Gamma\).  If a second DXZ
object has raw transform \(\phi\), the injectivity in
\Cref{thm:pseff-divisorial-inversion} identifies it with
\(\boldsymbol\Gamma\); this proves both uniqueness and independence of the
approximating model net.  Finally, for \(\phi,\psi\in
\PSH_{\mathrm{MP}}(\alpha)\), the equivalence
\eqref{eq:pseff-order-reflection} applied to their all-point raw transforms
gives
\[
 \mathcal S_\alpha(\phi)\leq\mathcal S_\alpha(\psi)
 \quad\Longleftrightarrow\quad
 \phi\leq\psi.
\]
Thus the section preserves and reflects order.
\end{proof}

\begin{theorem}[K\"ahler truncation closure]
\label{thm:p143-full-kahler-comparison}
In the setting of \Cref{thm:p143-mp-section}, the raw
ideal-valuative transform is an order isomorphism
\[
 \mathfrak T_\alpha:
 \PSH^{\NA}_{\mathrm{DXZ}}(X,\omega)
 \xrightarrow{\ \sim\ }
 \PSH_{\mathrm{MP}}(\alpha).
\tag{71.41j}\label{eq:p143-full-kahler-comparison}
\]
Its inverse is the section \(\mathcal S_\alpha\) constructed there.
More precisely, let \(\mathbf 0\) denote the constant DXZ object of value
zero and, for \(\boldsymbol\Gamma\in
\PSH^{\NA}_{\mathrm{DXZ}}(X,\omega)\), put
\[
 \boldsymbol\Gamma^{(j)}
 :=\boldsymbol\Gamma\vee(\mathbf 0-j),
 \qquad j\geq1.
\tag{71.41k}\label{eq:p143-dxz-truncation}
\]
Then every \(\boldsymbol\Gamma^{(j)}\) has finite energy,
\(\boldsymbol\Gamma^{(j)}\downarrow\boldsymbol\Gamma\), and
\[
 \mathfrak T_{\boldsymbol\Gamma}
 =\inf_{j\geq1}\mathfrak T_{\boldsymbol\Gamma^{(j)}}
 \quad\text{on }X^{\bethan},
\tag{71.41l}\label{eq:p143-raw-truncation-limit}
\]
where every function on the right belongs to
\(\mathcal E^1_{\mathrm{MP}}(\alpha)\).  On the finite-energy locus this
isomorphism is the correspondence of
\Cref{thm:pseff-kahler-finite-energy-link}, and the energies agree.
\end{theorem}

\begin{proof}
Fix \(\boldsymbol\Gamma\), write
\(c=\boldsymbol\Gamma_{\max}\), and set
\(c_j=\max\{c,-j\}\).  The constant object in
\eqref{eq:p143-dxz-truncation} is the one constructed in
\Cref{eq:pseff-constant-transform}.  In every K\"ahler perturbation class
\(\omega+\eta\), its level is zero for \(\tau<-j\).  The maximum
\(\boldsymbol\Gamma^{(j)}\) dominates this constant object.  Since its
\(\mathcal I\)-model levels are normalized to be at most zero, one has
\[
       \bigl(\boldsymbol\Gamma^{(j),\omega+\eta}\bigr)_\tau=0
       \qquad(\tau<-j).
\tag{71.41m}\label{eq:p143-truncation-zero-tail}
\]
Consequently its \(-\infty\)-level has full mass in every perturbed
component and
\[
             \vol(\boldsymbol\Gamma^{(j)})
             =V_\alpha:=\int_X\omega^n>0.
\]
The positive-volume recovery theorem
\Cref{thm:tna-positive-volume} therefore represents
\(\boldsymbol\Gamma^{(j)}\) by an unperturbed positive-mass test curve
\(\gamma^{(j)}\) in the class \(\omega\).  The recovery formula
\eqref{eq:tna-positive-recovery} and
\eqref{eq:p143-truncation-zero-tail} give
\(\gamma^{(j)}_\tau=0\) for \(\tau<-j\).  Hence
\[
 \mathbf E^0(\gamma^{(j)})
 =c_jV_\alpha+
   \int_{-j}^{c_j}
      \bigl(M_\omega(\gamma^{(j)}_\tau)-V_\alpha\bigr)\,\rd\tau
 \in[-jV_\alpha,c_jV_\alpha].
\tag{71.41n}\label{eq:p143-truncation-energy-bound}
\]
Thus \(\boldsymbol\Gamma^{(j)}\) has finite energy.  Notice that this
argument uses the actual zero tail of the recovered curve; it does not
infer finite energy from a merely divisorial lower bound.

Maximum is monotone, so the objects in
\eqref{eq:p143-dxz-truncation} decrease with \(j\).  Their endpoints are
\(c_j\downarrow c\), hence their decreasing infimum exists.  At every
positive rational divisorial point \(q\),
\Cref{prop:pseff-operation-formulas} gives
\[
 \begin{split}
 \mathfrak T_{\inf_j\boldsymbol\Gamma^{(j)}}(q)
   &=\inf_j\max\{\mathfrak T_{\boldsymbol\Gamma}(q),-j\}\\
   &=\mathfrak T_{\boldsymbol\Gamma}(q).
 \end{split}
\]
The endpoint formula gives the same equality at the trivial
semivaluation.  Injectivity in
\Cref{thm:pseff-divisorial-inversion} therefore yields
\[
                         \inf_j\boldsymbol\Gamma^{(j)}
                         =\boldsymbol\Gamma.
\tag{71.41o}\label{eq:p143-dxz-truncation-limit}
\]

Let \(\phi_j\in\mathcal E^1_{\mathrm{MP}}(\alpha)\) correspond to
\(\boldsymbol\Gamma^{(j)}\) under
\eqref{eq:pseff-finite-energy-bijection}.  By
\Cref{cor:p143-all-point-kahler-comparison},
\[
                    \phi_j
                    =\mathfrak T_{\boldsymbol\Gamma^{(j)}}
                    \quad\text{on }X^{\bethan}.
\]
The sequence \((\phi_j)\) is decreasing by monotonicity of the raw
transform.  Moreover,
\[
                    \inf_j\phi_j(0)=\inf_jc_j=c> -\infty,
\]
so \(\phi:=\inf_j\phi_j\) is not identically \(-\infty\).
\Cref{prop:mp-psh-decreasing-stability} gives
\(\phi\in\PSH_{\mathrm{MP}}(\alpha)\).  The divisorial maximum formula
and the endpoint formula show
\[
                    \phi=\mathfrak T_{\boldsymbol\Gamma}
                    \quad\text{on }X^{\divis}.
\]
By
\Cref{thm:p142-raw-transform-usc,cor:p142-pseudoeffective-raw-usc},
\(\mathfrak T_{\boldsymbol\Gamma}\) is upper semicontinuous.  It is
antitone for the ideal order by its defining multiplier-ideal formula.
The comparison clause of
\Cref{thm:p143-dominated-divisorial-density} now promotes the preceding
equality to all of \(X^{\bethan}\).  This proves
\eqref{eq:p143-raw-truncation-limit} and shows that the raw transform maps
the entire DXZ cone into the MP cone.

\Cref{thm:p143-mp-section} already gives
\(\mathfrak T_\alpha\circ\mathcal S_\alpha=\mathrm{id}\).  The raw
transform is injective by
\Cref{thm:pseff-divisorial-inversion}; hence
\(\mathcal S_\alpha\circ\mathfrak T_\alpha=\mathrm{id}\) as well.
Pointwise monotonicity of the raw transform and
\eqref{eq:pseff-order-reflection} show that both maps preserve and reflect
order.  Finally, on the finite-energy locus the divisorial identity
\eqref{eq:pseff-finite-energy-divisorial-link} and injectivity identify
this map with \Cref{thm:pseff-kahler-finite-energy-link}; its energy
identity therefore applies.
\end{proof}

\begin{remark}\label{rem:p143-full-kahler-scope}
The theorem settles the P14.3 image question in a K\"ahler class without
projectivity, rationality, or a finite-energy hypothesis.
It does not extend the relatively nef MP cone unchanged to a big non-nef
class, nor does it construct an unnormalized energy or Monge--Amp\`ere
operator for a zero-volume class.
\end{remark}

There is a complementary algebraic statement on the
\emph{DXZ--Boucksom--Jonsson} side, including non-nef classes.  Assume that $X$ is smooth
projective and that
$\alpha=c_1(L)$ for a pseudoeffective $\mathbf Q$-line bundle $L$;
neither ampleness nor positive volume is required.  Then
\cite[Theorem~3.14]{DXZ23}, restated as \Cref{thm:comp-dxz}, is an order
isomorphism
\[
 \PSH^{\NA}_{\mathrm{DXZ}}(X,\theta)
 \xrightarrow{\ \sim\ }
 \PSH(L^{\mathrm{an,triv}}).
\tag{71.41}\label{eq:pseff-projective-dxz-bj-bijection}
\]
Here the raw function is also Boucksom--Jonsson psh at every
semivaluation. Its upper semicontinuity in arbitrary pseudoeffective
classes was proved in \Cref{cor:p142-pseudoeffective-raw-usc}; the new
input in this projective comparison is positivity. GAGA and
\Cref{thm:ts-algebraic-recovery} identify
$X^{\bethan}$ with $X^{\mathrm{an,triv}}$ and identify every
coherent-ideal evaluation.  Under this identification,
\cite[formula~(3.8) and Theorem~3.12]{Xia23Operations} defines the
Boucksom--Jonsson metric by the same formula as
\eqref{eq:pseff-inverse-system-transform}.  Consequently
$\mathfrak T_{\boldsymbol\Gamma}$ is upper semicontinuous and
Boucksom--Jonsson psh on all of $X^{\bethan}$, and
\eqref{eq:pseff-projective-dxz-bj-bijection} is literally the map
$\boldsymbol\Gamma\mapsto\mathfrak T_{\boldsymbol\Gamma}$, not merely
an equality on divisorial points.

The directed-limit argument of
\cite[Remark~3.15]{DXZ23} extends this identification from
$\mathbf Q$-line bundles to every pseudoeffective real
$(1,1)$-class in the N\'eron--Severi space of a smooth projective
manifold, with target
$\PSH(\alpha^{\NA})$.  This is the strongest proved projective
framework.  At this real-class level we assert the canonical
directed-limit identification stated in the source; the literal
multiplier-ideal formula on every semivaluation is claimed above only in
the $\mathbf Q$-line-bundle case.  The result does not cover a genuinely
transcendental Bott--Chern class outside the real N\'eron--Severi space.

The next theorem supplies the missing comparison in the natural ample
projective numerical framework.  It is a source-explicit synthesis of the
model descriptions and approximation theorems in
\cite{MP24,BJ22,DXZ23}; no novelty or priority claim is attached to it.

\begin{theorem}\label{thm:pseff-projective-full-mp-bj}
Let $X$ be a connected smooth projective variety over $\mathbf C$, and
let
\[
                 \alpha\in\operatorname{Amp}(X)
                         \subset N^1(X)_{\mathbf R}.
\]
On the Mesquita--Piccione side, view $\alpha$ as its real Bott--Chern
class.  Give $\mathbf C$ the trivial absolute value, and let
\[
 \iota:X^{\bethan}\xrightarrow{\ \sim\ }X^{\mathrm{an,triv}}
\]
be the canonical homeomorphism induced by GAGA.  Write
$\PSH_{\mathrm{MP}}(\alpha)$ for the decreasing completion of rational
PL functions whose cohomological model class is relatively nef, and write
$\PSH_{\mathrm{BJ}}(\alpha):=\PSH(\alpha^{\NA})$ for the
Boucksom--Jonsson cone.  Then pullback by $\iota$ is an order
isomorphism
\[
 \iota^*\PSH_{\mathrm{BJ}}(\alpha)
       =\PSH_{\mathrm{MP}}(\alpha).
\tag{71.42}\label{eq:pseff-projective-full-mp-bj}
\]
More precisely, it first identifies the rational PL model cones:
\[
 \iota^*\bigl(\PL(X^{\mathrm{an,triv}})
          \cap\PSH_{\mathrm{BJ}}(\alpha)\bigr)
 =\PL(X^{\bethan})\cap\PSH_{\mathrm{MP}}(\alpha),
\tag{71.43}\label{eq:pseff-projective-model-mp-bj}
\]
and the two sides of
\eqref{eq:pseff-projective-full-mp-bj} are the same decreasing completion
of this common rational PL cone.

If $L$ is an ample $\mathbf Q$-line bundle with
$c_1(L)=\alpha$, and $\theta$ is a smooth representative of
$c_1(L)$, then composition with the unrestricted
Darvas--Xia--Zhang--Boucksom--Jonsson comparison gives an order
isomorphism
\[
 \PSH^{\NA}_{\mathrm{DXZ}}(X,\theta)
       \xrightarrow{\ \sim\ }
 \PSH_{\mathrm{MP}}(c_1(L))
\tag{71.44}\label{eq:pseff-projective-dxz-mp-bijection}
\]
on all semivaluations, with no finite-energy hypothesis.
\end{theorem}

\begin{proof}
We first record the source facts used in the comparison.  Mesquita--Piccione
identify rational PL functions on $X^{\bethan}$ with the direct limit of
invariant vertical $\mathbf Q$-Cartier divisors on dominating analytic
test configurations
\cite[Proposition~2.4.3 and Theorem~2.6.4]{MP24}.  Their
Definition~4.1.2 declares such a function $\alpha$-psh precisely when
the class $\alpha_{\mathcal X}+[D]$ is relatively nef, and their
Definition~4.3.1 defines the full cone by pointwise decreasing nets in
that rational PL model cone, excluding the identically $-\infty$ limit.

Boucksom--Jonsson define PL positivity by the formally analogous algebraic
relative-nef condition on algebraic test configurations
\cite[Definition~3.2]{BJ22}.  Their Lemma~4.3 identifies this PL condition
with membership in their full psh cone.  Their Theorem~4.5 gives closure
under decreasing nets whose limit is generically finite.  Finally, when
$\alpha$ is ample, their Theorem~4.15(iii) writes every
$\alpha$-psh function as the pointwise decreasing limit of a net in
$\mathcal H^{\mathrm{dom}}(\alpha)$.  By Definition~3.8 in that source,
each member of $\mathcal H^{\mathrm{dom}}(\alpha)$ is a rational PL
function represented on a dominating algebraic test configuration by a
vertical $\mathbf Q$-Cartier divisor $D$ for which
$\alpha_{\mathcal X}+D$ is relatively ample.

We compare the rational PL lattices and their positive cones.  The
Mesquita--Piccione rational PL functions form the $\mathbf Q$-vector
space spanned by flag-ideal functions.  Choose the finitely many flag
ideals occurring in one presentation and clear all denominators.  Each
flag ideal on $X\times\mathbf P^1$ has an expansion
\[
                    \mathfrak a=\sum_j\mathfrak a_jt^j
\]
with only finitely many nontrivial coherent coefficient ideals on $X$.
Since $X$ is projective, GAGA algebraizes each $\mathfrak a_j$, and
the finite weighted expansion algebraizes the flag ideal together with
its $\mathbf C^*$-action.  Take one projective model principalizing all
the resulting algebraic flag ideals.  Blowup is the relative Proj of the
Rees algebra, so analytification commutes with this blowup; proper GAGA for
finite algebras likewise identifies algebraic and analytic normalizations.
A $\mathbf C^*$-equivariant projective resolution in characteristic zero
then gives a common smooth projective algebraic test configuration carrying
the pulled-back vertical divisors and their rational linear combination.
Conversely, analytifying an algebraic flag ideal gives the corresponding
analytic flag ideal.  Restriction from $\mathbf P^1$ to
$\mathbf A^1$ gives a Boucksom--Jonsson test configuration, while the
equivariantly trivial compactification over $\mathbf P^1$ reverses the
passage; all vertical divisors in question are supported over $0$.
The coefficient formula on vertical components,
\[
                    D=\sum_E b_E\varphi(r(v_E))E,
\]
is the same in both theories.  Hence $\iota$ identifies their rational
PL lattices, and every rational PL function may be tested on one common
smooth projective algebraic model.  This only algebraizes flag-ideal
presentations; it does not assert that every analytic test configuration
is algebraizable.  The underlying space and ideal evaluations agree by
\cite[Remark~1.3.3]{MP24}.

Let $\mathcal Y\to X\times\mathbf P^1$ be such a common smooth
projective model, let $D$ represent $\varphi$, and put
\[
                         \beta=\alpha_{\mathcal Y}+[D].
\]
We next justify checking MP positivity on this algebraic determination.
Suppose an analytic determination
$(\mathcal X,D_{\mathcal X})$ originally represents the same
$\varphi$.  Cofinality and directedness of projective snc analytic test
configurations \cite[Remark~2.5.2]{MP24} give a common analytic domination
\[
 \mathcal Z\xrightarrow{p}\mathcal X,
 \qquad
 \mathcal Z\xrightarrow{q}\mathcal Y.
\]
After one further refinement if necessary, injectivity in the direct-limit
isomorphism \cite[Theorem~2.6.4]{MP24} gives
\[
                  p^*D_{\mathcal X}=q^*D.
\]
The two pulled-back cohomological classes on $\mathcal Z$ therefore
agree.  Relative nefness ascends under pullback and descends under these
proper modifications; this is the model-independence assertion in
\cite[Definition~4.1.2 and the proof of Theorem~4.1.5]{MP24}.
Consequently the MP class on $\mathcal X$ is relatively nef if and only
if the analytified class $\beta$ on $\mathcal Y$ is relatively nef.

On a smooth projective variety, the analytic and algebraic nef cones agree
on $N^1_{\mathbf R}$.  Indeed, algebraic ample classes are K\"ahler by
Kodaira.  Conversely, let $\gamma\in N^1_{\mathbf R}$ be K\"ahler and
fix an ample class $H$.  Openness of the K\"ahler cone gives
$\varepsilon>0$ such that $\gamma-\varepsilon H$ is still K\"ahler.
Thus $\gamma\cdot C\geq\varepsilon H\cdot C$ for every irreducible
curve $C$.  Positive homogeneity and continuity extend this inequality
to the closed cone of curves, and Kleiman's criterion places $\gamma$
in the algebraic ample cone.  Thus the K\"ahler and ample cones agree
inside $N^1_{\mathbf R}$; taking closures identifies analytic and
algebraic nefness.

Apply this observation to every irreducible component of the
simple-normal-crossings central fibre of $\mathcal Y$.  Each component
is smooth and projective.  Restriction of the compactified model to
$\mathbf A^1$ gives the Boucksom--Jonsson test configuration, while the
canonical equivariantly trivial compactification recovers the MP model.
On every noncentral fibre, including the fibre over infinity,
$\beta$ restricts to the ample class $\alpha$.  Hence $\beta$ is
relatively nef in the analytic sense used by Mesquita--Piccione if and
only if it is relatively nef in the algebraic sense used by
Boucksom--Jonsson; compare the algebraic central-fibre criterion
\cite[Lemma~3.1]{BJ22}.  The common-refinement argument applies this
criterion even when the original analytic representative of $\varphi$
was carried by a nonalgebraizable test configuration.  This proves
\eqref{eq:pseff-projective-model-mp-bj}.

We pass to decreasing completions.  Suppose first that
$\varphi\in\PSH_{\mathrm{MP}}(\alpha)$.  Choose a net
$\varphi_i\downarrow\varphi$ in the MP rational PL model cone.  By
\eqref{eq:pseff-projective-model-mp-bj} and
\cite[Lemma~4.3]{BJ22}, every $\varphi_i$ is BJ
$\alpha$-psh.  The limit is upper semicontinuous because it is a
decreasing limit of continuous functions.  It is also generically finite.
Since $X$ is connected and smooth, it is irreducible, and the PL maximum
principle \cite[Proposition~3.6(iv)]{BJ22} gives
\[
                  \varphi_i(v)\leq\varphi_i(v_{\mathrm{triv}})
                  \qquad(v\in X^{\mathrm{an,triv}}).
\]
If $\varphi(v_{\mathrm{triv}})=-\infty$, passage to the limit would give
$\varphi\equiv-\infty$, contrary to the MP definition.  The limit is
therefore finite at the trivial valuation and hence generically finite by
\cite[Lemma~1.4(ii)]{BJ22}.  The decreasing-closure theorem
\cite[Theorem~4.5]{BJ22} yields
$\varphi\in\PSH_{\mathrm{BJ}}(\alpha)$.

Conversely, let $\psi\in\PSH_{\mathrm{BJ}}(\alpha)$.  By
\cite[Theorem~4.15(iii)]{BJ22}, there is a decreasing net
\[
                 \psi_i\in\mathcal H^{\mathrm{dom}}(\alpha),
                 \qquad \psi_i\downarrow\psi.
\]
Each $\psi_i$ is rational PL and relatively ample on an algebraic
dominating model.  In particular it lies in the common relatively nef PL
cone of \eqref{eq:pseff-projective-model-mp-bj}.  Since a BJ psh function
is generically finite, $\psi\not\equiv-\infty$.  The defining
decreasing-net criterion of \cite[Definition~4.3.1]{MP24} gives
$\psi\in\PSH_{\mathrm{MP}}(\alpha)$.  This proves
\eqref{eq:pseff-projective-full-mp-bj}; the identification visibly
preserves and reflects pointwise order.

There is no conflict with \cite[Remark~4.1.3]{MP24}, which says that the MP
relatively K\"ahler model locus can be larger than the algebraic
Fubini--Study locus.  The proof uses equality only of the rational
relatively nef PL cones and the cofinal BJ regularization by
$\mathcal H^{\mathrm{dom}}(\alpha)$; it does not identify those two
smaller strictly positive loci.

Finally, when $\alpha=c_1(L)$ for an ample $\mathbf Q$-line bundle
and $\theta$ is a smooth representative, \cite[Theorem~3.14]{DXZ23}
identifies the full DXZ cone with the full Boucksom--Jonsson cone on
$X^{\mathrm{an,triv}}$, without an energy hypothesis.  Composing this
order isomorphism with \eqref{eq:pseff-projective-full-mp-bj} proves
\eqref{eq:pseff-projective-dxz-mp-bijection}.
\end{proof}

The projective comparison still does not create an MP cone in a non-nef
class; \Cref{prop:pseff-relative-nef-obstruction} rules out the unchanged
relative-nef model signature there.  The exact comparison ladder is now:
\begin{enumerate}
\item for a compact K\"ahler pseudoeffective class, the transform is
      canonical on all semivaluations and determining on the divisorial
      locus;
\item for a projective pseudoeffective real N\'eron--Severi class, the full
      DXZ space is canonically identified with the Boucksom--Jonsson psh
      space; for a $\mathbf Q$-line-bundle class this identification is
      literally the ideal formula
      \eqref{eq:pseff-inverse-system-transform} on every semivaluation;
\item for every K\"ahler class, the raw transform identifies the full
      DXZ and MP cones on all semivaluations by
      \Cref{thm:p143-full-kahler-comparison}; on the finite-energy locus,
      the class-normalized radial and non-Archimedean energies agree;
\item for a smooth projective real ample numerical class, the full MP and
      Boucksom--Jonsson cones agree; for an ample $\mathbf Q$-class, the
      full DXZ--MP identification
      \eqref{eq:pseff-projective-dxz-mp-bijection} requires no energy
      hypothesis.
\end{enumerate}

\section{Operations on the verified domain}
\label{sec:pseff-operations}

The ideal-valuative transform makes translation transparent on every
semivaluation.  The other pointwise formulas proved in
\cite{Xia23Operations} are divisorial formulas; we retain that domain.

\begin{proposition}\label{prop:pseff-operation-formulas}
Let \(\boldsymbol\Gamma,\boldsymbol\Gamma'\in
\PSH^{\NA}(X,\theta)\), let
\(\boldsymbol\Lambda\in\PSH^{\NA}(X,\theta')\), and let \(C\in\mathbf R\).
Then
\[
 \mathfrak T_{\boldsymbol\Gamma+C}(v)
 =\mathfrak T_{\boldsymbol\Gamma}(v)+C
 \qquad(v\in X^{\bethan}).
\tag{71.45}\label{eq:pseff-operation-translation}
\]
For every rational divisorial \(v\),
\[
 \begin{aligned}
 \mathfrak T_{\boldsymbol\Gamma\vee\boldsymbol\Gamma'}(v)
   &=\max\{\mathfrak T_{\boldsymbol\Gamma}(v),
            \mathfrak T_{\boldsymbol\Gamma'}(v)\},\\
 \mathfrak T_{\boldsymbol\Gamma+\boldsymbol\Lambda}(v)
   &=\mathfrak T_{\boldsymbol\Gamma}(v)
     +\mathfrak T_{\boldsymbol\Lambda}(v),\\
 \mathfrak T_{\lambda\boldsymbol\Gamma}(v)
   &=\lambda\mathfrak T_{\boldsymbol\Gamma}(v),
 \qquad \lambda>0.
 \end{aligned}
\tag{71.46}\label{eq:pseff-operation-algebra}
\]
If \((\boldsymbol\Gamma^i)_{i\in I}\) is nonempty and
\(\sup_i\boldsymbol\Gamma^i_{\max}<+\infty\), then
\[
 \mathfrak T_{\uscsup_i\boldsymbol\Gamma^i}(v)
 =\sup_i\mathfrak T_{\boldsymbol\Gamma^i}(v)
\tag{71.47}\label{eq:pseff-operation-supremum}
\]
at every rational divisorial \(v\).  If
\((\boldsymbol\Gamma^i)_{i\in I}\) is a decreasing net and
\(\inf_i\boldsymbol\Gamma^i_{\max}>-\infty\), then
\[
 \mathfrak T_{\inf_i\boldsymbol\Gamma^i}(v)
 =\inf_i\mathfrak T_{\boldsymbol\Gamma^i}(v)
\tag{71.48}\label{eq:pseff-operation-infimum}
\]
at every rational divisorial \(v\).
\end{proposition}

\begin{proof}
Equation \eqref{eq:pseff-operation-translation} was proved directly from
multiplier ideals in \Cref{prop:pseff-one-curve-transform} and survives
componentwise passage to the inverse system by
\Cref{thm:pseff-perturbation-independent-transform}.

At a divisorial point, \eqref{eq:pseff-inverse-system-transform} is the
associated non-Archimedean potential of
\cite[formula~(3.7)]{Xia23Operations}.  The maximum formula is
\cite[Corollary~4.8]{Xia23Operations}; the addition and rescaling formulas
are \cite[Lemma~4.17]{Xia23Operations}.  These results apply to the
componentwise operations of
\Cref{def:operations-addition,def:operations-maximum} and give
\eqref{eq:pseff-operation-algebra}.

Under the displayed upper endpoint bound, the regularized supremum exists
by \Cref{def:operations-supremum}, and
\cite[Theorem~4.13]{Xia23Operations} gives
\eqref{eq:pseff-operation-supremum}.  Under the displayed lower endpoint
bound, the decreasing infimum exists by
\Cref{def:operations-infimum}, and
\cite[Lemma~4.10]{Xia23Operations} gives
\eqref{eq:pseff-operation-infimum}.  All cited identities hold for every
positive real multiple of a divisorial valuation; restricting to rational
multiples yields the statement.  None of these citations provides the
same raw pointwise formulas at arbitrary nondivisorial semivaluations.
\end{proof}

\section{Raw filtration curves in a big class}
\label{sec:pseff-filtration-maximality}

We use squared Hermitian norms and
\(\mathrm d\mathrm d^{\mathrm c}
 =\frac{\mathrm i}{2\pi}\partial\bar\partial\).
Let \(X\) be a connected smooth complex projective variety, let \(L\) be a
big line bundle, let \(h\) be a smooth Hermitian metric, and put
\(\theta=c_1(L,h)\).  Set
\[
 V_\theta=\sup\{v\in\operatorname{PSH}(X,\theta):v\leq0\}.
\]
Write \(v\preceq u\) if \(v\leq u+C\) for some real constant \(C\), and use
the normalized singularity envelope
\[
 P_\theta[u]
 =\left(\sup\{v\in\operatorname{PSH}(X,\theta):
                  v\leq0,\ v\preceq u\}\right)^*.
\]
Here \({}^*\) denotes upper semicontinuous regularization.  The obstacle
\(0\) gives the same envelope as \(V_\theta\), since every
\(\theta\)-psh function below \(0\) is below \(V_\theta\).
The multiplier-ideal envelope \(P_\theta[u]_{\mathcal I}\) is the
normalized envelope with the conditions
\(\mathcal I(tv)\subseteq\mathcal I(tu)\) for every \(t\geq0\).

Let \(\mathcal F\) be a decreasing, left-continuous, multiplicative,
linearly bounded filtration of \(R(X,L)\), precisely as in
\cite[Definition~4.1]{DXZ23}.  Put
\[
 c=\tau_L(\mathcal F),\qquad
 W_m^\tau=\mathcal F^{m\tau}H^0(X,L^{\otimes m}),
\]
and, for \(\tau<c\), use the normalization of
\cite[Definition~4.3]{DXZ23}:
\[
\tag{71.48a}\label{eq:pseff-filtration-raw-definition}
 u_{\tau,m}
 =\sup_{\substack{s\in W_m^\tau\\
             \sup_X|s|_{h^m}^2\leq1}}
       \frac1m\log|s|_{h^m}^2,\qquad
 u_\tau=\left(\sup_{m\geq1}u_{\tau,m}\right)^*.
\]
An empty nonzero section space gives \(-\infty\).
For a nonzero finite-dimensional space the first supremum gives a
\(\theta\)-psh function with an analytic singularity type.  Notice that the
factor in \eqref{eq:pseff-filtration-raw-definition} is \(1/m\), not
\(1/(2m)\).

\begin{theorem}\label{thm:pseff-filtration-maximality}
With these hypotheses, for every \(\tau<c\),
\[
\tag{71.48b}\label{eq:pseff-filtration-maximality}
 u_\tau=P_\theta[u_\tau]=P_\theta[u_\tau]_{\mathcal I}.
\]
In particular, the raw filtration test curve is maximal throughout its
strict interior, which is the positive-mass range.
\end{theorem}

The proof establishes the ordinary equality
\(u_\tau=P_\theta[u_\tau]\).  The second equality is
\cite[Lemma~4.5]{DXZ23}.  No finite-generation hypothesis on
\(\mathcal F\) is needed.

\begin{proof}

The new point is a function-level estimate for a fixed birational linear
system.  It uses the actual sup norms, not just the dimension or
multiplier ideals of the section spaces.

If \(\dim X=0\), connectedness makes \(X\) a point.  Every nonzero
section space has normalized evaluation supremum \(0\), so every strict
interior slice and its two envelopes equal \(0\).  The terminal assertion
is immediate as well.  Henceforth assume \(\dim X\geq1\).

\paragraph{The fixed-degree estimate.}

Let \(W=\bigoplus_{a\geq0}W_a\subseteq R(X,L)\) be a multiplicative graded
linear series.  For every \(a\geq1\), set
\[
 u_a=\sup_{\substack{s\in W_a\\\|s\|_{h^a,\infty}\leq1}}
                   a^{-1}\log|s|_{h^a}^2,\qquad
 u_W=(\sup_{a\geq1}u_a)^*.
\]
Assume \(u_W\not\equiv-\infty\).  Each \(u_a\leq0\), hence \(u_W\leq0\).
Fix a positive integer \(m\) for which the rational map defined by \(W_m\)
is birational onto its image.  We will prove
\[
\tag{71.48c}\label{eq:pseff-filtration-fixed-degree-bound}
                         P_\theta[u_m]\leq u_W.
\]

\paragraph{Step 1: resolve the fixed linear system and its obstacle.}
Take a projective resolution of the base ideal of \(W_m\),
\(\mu:Y\to X\), with \(Y\) smooth and connected.  There are an effective
Cartier divisor \(E\), a globally generated line bundle \(A\), and an
isomorphism
\[
\tag{71.48d}\label{eq:pseff-filtration-resolved-bundle}
             \mu^*L^{\otimes m}=A\otimes\mathcal O_Y(E).
\]
Dividing the pulled-back sections by the canonical section \(s_E\)
gives a basepoint-free subspace \(V\subseteq H^0(Y,A)\).
Let \(f:Y\to Z\subseteq\mathbf P(W_m^*)\) be its morphism.
Then \(f\) is birational, \(Z\) is integral, and
\(A=f^*\mathcal O_Z(1)\).  In particular \(A\) is big as well as globally
generated.  We do not assume that \(Z\) is normal or that \(V\) is the
complete linear system.

Give \(A\) a smooth Fubini--Study metric \(a\) from \(V\);
\(\alpha=c_1(A,a)\) is semipositive.
Use \eqref{eq:pseff-filtration-resolved-bundle} to define the smooth metric \(h_E\) on
\(\mathcal O_Y(E)\) by \(\mu^*h^m=a\otimes h_E\), and put
\[
 g=\log|s_E|_{h_E}^2,\qquad \chi=-g.
\]
Thus \(\chi\) is a lower semicontinuous function, smooth off \(E\),
bounded below, and equal to \(+\infty\) on \(E\).
Moreover
\[
 \alpha+c_1(\mathcal O_Y(E),h_E)=m\mu^*\theta,\qquad
 \mathrm d\mathrm d^{\mathrm c}g=[E]-c_1(\mathcal O_Y(E),h_E).
\]
There is a continuous bounded function \(v_m\in\operatorname{PSH}(Y,\alpha)\)
such that
\[
\tag{71.48e}\label{eq:pseff-filtration-fixed-singularity}
                       m\mu^*u_m=g+v_m.
\]
Indeed, \(v_m\) is the logarithm of the dual norm of evaluation on the
finite-dimensional normed space \(V\).  Evaluation never vanishes on
\(Y\); compactness gives both bounds.  Equivalence with a Euclidean norm
also shows directly the stated analytic singularity type.

Define
\[
\tag{71.48f}\label{eq:pseff-filtration-residual-envelope}
 Q=\left(\sup\{q\in\operatorname{PSH}(Y,\alpha):q\leq\chi\}\right)^*.
\]
This is a bounded \(\alpha\)-psh function.  For the lower bound, use a
constant below \(\inf_Y\chi\), since \(\alpha\geq0\).
For the upper bound, choose a nonempty relatively compact coordinate ball
outside \(E\), on which \(\chi\) is bounded above.
The usual \(L^1\)-compactness of the normalized \(\alpha\)-psh functions
implies a uniform global upper bound for all \(q\leq\chi\): otherwise,
subtracting unbounded maxima would give normalized functions tending to
\(-\infty\) on that ball, contradicting this compactness.
Consequently the supremum is locally uniformly bounded above.
Its regularization satisfies \(Q\leq\chi\) off \(E\), where \(\chi\) is
continuous; on \(E\) the inequality is automatic.  Thus \(Q\) itself is
an admissible competitor.

The definitions give the following exact identity of psh potentials:
\[
\tag{71.48g}\label{eq:pseff-filtration-envelope-identification}
                   \mu^*P_\theta[u_m]=\frac{g+Q}{m}.
\]
Here are the details, including descent.
If \(v\leq0\) is \(\theta\)-psh and \(v\preceq u_m\), then
\(m\mu^*v-g\) is \(\alpha\)-psh away from \(E\).
By \eqref{eq:pseff-filtration-fixed-singularity} it is locally bounded above near
\(E\), so its local psh weights extend across \(E\).
The extension is at most \(\chi\), so is at most \(Q\).
Conversely, for any \(\alpha\)-psh \(q\leq\chi\), the potential
\((g+q)/m\) is \(\mu^*\theta\)-psh and is at most \(0\).
It descends to a \(\theta\)-psh function on \(X\).
To see the last assertion without a smoothness assumption on an image
linear system, work on the open set where \(\mu\) is an isomorphism.
The descended local psh weights are bounded above near the exceptional
image, because their pullbacks are bounded above over compact
neighbourhoods.  Removable singularities extend them across that analytic
set.  Pulling back these extensions recovers the original weights:
both are psh and agree outside a proper analytic subset.
Finally \(q\) is bounded above and \(v_m\) is bounded below, so the
descended function is \(\preceq u_m\).
The two constructions identify the competitors.
Their upper envelopes give \eqref{eq:pseff-filtration-envelope-identification}
first away from the exceptional and divisor sets, then everywhere by
equality of psh functions almost everywhere.

\paragraph{Step 2: quantize this extended-real obstacle.}
For \(k\geq1\), define the genuine norm
\[
 \|s\|_{\chi,k}
   =\sup_Y |s|_{a^k}\exp(kg/2),\qquad s\in H^0(Y,A^{\otimes k}),
\]
and the finite-degree potential
\[
 b_{\chi,k}
  =\sup_{\|s\|_{\chi,k}\leq1}k^{-1}\log|s|_{a^k}^2.
\]
The norm is finite since \(\exp(g)\) is smooth and bounded.
It is definite since it is positive on a nonempty open set and a
holomorphic section vanishing there vanishes identically.
Every \(b_{\chi,k}\) is at most \(Q\).

We claim that
\[
\tag{71.48h}\label{eq:pseff-filtration-extended-obstacle-quantization}
                         b_{\chi,k}\longrightarrow Q
                  \quad\hbox{almost everywhere on }Y.
\]
This only needs the smooth-obstacle Bergman theorem for a big line bundle.
For clarity, set
\[
             \chi_j=-\log\bigl(|s_E|_{h_E}^2+j^{-1}\bigr).
\]
These smooth functions increase to \(\chi\).
Let \(Q_j=P_\alpha(\chi_j)\).
They increase to \(Q\) after upper semicontinuous regularization.
Indeed, for any \(q\leq\chi\) and \(\varepsilon>0\), the increasing open
sets
\(\{q-\varepsilon<\chi_j\}\) cover the compact space \(Y\).
One of them is all of \(Y\), so \(q-\varepsilon\leq Q_j\).
Taking suprema and then regularizations proves the assertion.
In particular \(Q_j\to Q\) almost everywhere.

For each fixed \(j\), use the smooth metric \(a\exp(-\chi_j)\) on the big
line bundle \(A\).
The Bergman-weight convergence in \cite[Theorem~1.4]{Ber11} gives
the full-section \(L^2\) quantization of \(Q_j\), uniformly on compact
subsets of \(Y\setminus\mathbf B_+(A)\).
The auxiliary fixed bundle in that theorem is the trivial bundle.
There is no positivity assumption on this smooth metric.
The same assertion holds for the sup norm: with a fixed smooth volume
form of total mass one,
\[
 \|s\|_{2,j}\leq\|s\|_{\infty,j}
       \leq C_j k^n\|s\|_{2,j}
\]
for all \(k\) sufficiently large.  For the upper estimate cover \(Y\)
by coordinate charts with smooth local weights, apply the holomorphic
mean-value inequality on balls of radius \(k^{-1}\), and note that
\(k\) times the oscillation of a smooth weight on such a ball is bounded.
This gives a polynomial loss, sufficient here.
Consequently the two logarithmic evaluation suprema differ by at most
\(2\log(C_jk^n)/k\), which tends to zero.
Thus, if \(b_{\chi_j,k}\) uses the sup norm associated with \(\chi_j\),
\[
                  b_{\chi_j,k}\longrightarrow Q_j
             \quad\hbox{on }Y\setminus\mathbf B_+(A).
\]
Since \(\chi_j\leq\chi\), their unit balls satisfy
\(b_{\chi_j,k}\leq b_{\chi,k}\leq Q\).
Taking a lower limit in \(k\), then the supremum over the countable set
of \(j\), proves \eqref{eq:pseff-filtration-extended-obstacle-quantization}.
The exceptional set is the augmented base locus and the null set where
the increasing \(Q_j\)'s differ from their regularized limit.
Neither uniform convergence at \(E\) nor continuity of \(Q\) is used.

\paragraph{Step 3: return complete sections to the given graded series.}
This is where birationality matters.
Let \(\nu:Z^\nu\to Z\) be the normalization.  It is finite.
As \(Y\) is normal and \(f\) is proper birational,
\[
                         f_*\mathcal O_Y=\nu_*\mathcal O_{Z^\nu}.
\]
For example, factor \(f\) through \(Z^\nu\).
A function regular upstairs is a rational function on the normal target
without a pole at any prime divisor, hence is regular; this gives the
displayed identity.  Alternatively apply \cite[Tag~0AY8]{StacksProject}:
the morphism \(Y\to Z^\nu\) is proper, the target is normal and integral,
the source is integral and reduced, and its generic fibre is the single
point with function field \(\mathbf C(Z^\nu)\).

The coherent quotient
\(\nu_*\mathcal O_{Z^\nu}/\mathcal O_Z\) has zero generic stalk.
Its annihilator \(\mathfrak c\) is therefore a nonzero coherent ideal and
\[
              \mathfrak c\,\nu_*\mathcal O_{Z^\nu}\subseteq\mathcal O_Z.
\]
Since \(\mathcal O_Z(1)\) is ample, choose a fixed integer \(r\geq1\) and
a nonzero section
\[
                  b\in H^0(Z,\mathfrak c\otimes\mathcal O_Z(r)).
\]
Global generation of a sufficiently high twist guarantees its existence.
The projection formula gives, for every \(k\),
\[
\tag{71.48i}\label{eq:pseff-filtration-conductor-inclusion}
 b\,H^0(Y,A^{\otimes k})
                    \subseteq H^0(Z,\mathcal O_Z(k+r)),
\]
where pullback to \(Y\) is understood on the left.
For \(k\) large, Serre vanishing applied to the ideal of
\(Z\subseteq\mathbf P(W_m^*)\) makes the restriction of homogeneous
polynomials of degree \(k+r\) onto the right side surjective.
Therefore, for every \(s\in H^0(Y,A^{\otimes k})\),
\[
\tag{71.48j}\label{eq:pseff-filtration-polynomial-return}
                 s_E^{\,k+r}bs
          \ \in\ \mu^*\operatorname{im}
                    \bigl(\operatorname{Sym}^{k+r}W_m
                                      \longrightarrow W_{m(k+r)}\bigr).
\]
This inclusion has the indicated direction.  An arbitrary section of the
normalization does not descend without the conductor factor.
The standard algebraic inputs here are finite normalization and Serre
vanishing/global generation; see \cite[Tags~0BXQ, 01PR, and 02O1]{StacksProject}.
No finite generation of the entire \(W\) is involved.

Regard \(b\) now as its pullback, a nonzero section of \(A^{\otimes r}\), and
set
\[
 B=\sup_Y |b|_{a^r}\exp(rg/2),\qquad
 \beta=\log|b|_{a^r}^2.
\]
The constant satisfies \(0<B<\infty\).
If \(\|s\|_{\chi,k}\leq1\), the section in
\eqref{eq:pseff-filtration-polynomial-return} has \(\mu^*h^{m(k+r)}\)-sup norm at most
\(B\), since its pointwise norm is exactly
\[
       \bigl(|b|_{a^r}\exp(rg/2)\bigr)
       \bigl(|s|_{a^k}\exp(kg/2)\bigr).
\]
After division by \(B\) it is an allowed section in the definition of
\(u_W\).  Taking the supremum over \(s\) yields, away from \(E\cup\{b=0\}\),
\[
\tag{71.48k}\label{eq:pseff-filtration-pointwise-return}
 \mu^*u_W\geq
       \frac{g}{m}
       +\frac{k b_{\chi,k}+\beta-2\log B}{m(k+r)} .
\]
All degrees, norms, and filtration memberships in this inequality are
exact.  Only its limit discards the fixed cost \(r\) and the fixed
logarithmic term.
By \eqref{eq:pseff-filtration-extended-obstacle-quantization}, its right side converges
almost everywhere to \((g+Q)/m\).
Using \eqref{eq:pseff-filtration-envelope-identification} gives
\(u_W\geq P_\theta[u_m]\) almost everywhere on \(X\).
Both functions are \(\theta\)-psh, so the inequality holds everywhere.
This proves \eqref{eq:pseff-filtration-fixed-degree-bound}.

\paragraph{Application to the filtration.}

Fix \(\tau<c\) and take \(W_a=W_a^\tau\).
The series contains an ample series; consequently there is an integer
\(m_0\geq1\) such that the maps defined by \(W_{m_0q}^\tau\) are birational
for all sufficiently large \(q\).
This is the precise input stated in the proof of
\cite[Proposition~4.6]{DXZ23}, referring to \cite[Lemma~1.6]{BC11}.
For \(m=m_0 2^j\) sufficiently large,
\eqref{eq:pseff-filtration-fixed-degree-bound} gives
\[
                    P_\theta[u_{\tau,m_0 2^j}]\leq u_\tau.
\]
Multiplicativity and the sup-norm constraint give
\(u_{\tau,2^j}\leq u_{\tau,m_0 2^j}\): take the \(m_0\)-th power of any
normalized section.  Monotonicity of the singularity envelope therefore
gives
\[
\tag{71.48l}\label{eq:pseff-filtration-dyadic-bound}
                         P_\theta[u_{\tau,2^j}]\leq u_\tau.
\]
The proof of \cite[Lemma~4.5]{DXZ23} supplies the exact identity
\[
 P_\theta[u_\tau]
      =\left(\lim_{j\to\infty}P_\theta[u_{\tau,2^j}]\right)^*
      =P_\theta[u_\tau]_{\mathcal I}.
\]
In that proof the dyadic potentials increase almost everywhere to
\(u_\tau\); the maximized finite-level potentials are \(\mathcal I\)-model,
and positive mass makes the increasing-limit closure theorem applicable.
It is not an assertion that the raw potentials were already maximal.
Taking the upper semicontinuous regularization in
\eqref{eq:pseff-filtration-dyadic-bound} gives \(P_\theta[u_\tau]\leq u_\tau\).
The reverse inequality follows from \(u_\tau\leq0\), so
\eqref{eq:pseff-filtration-maximality} follows.

For completeness, positive mass on this whole strict interior is not an
extra restriction.  Choose \(a<\tau<b<c\) with \(u_a=V_\theta\), using
linear boundedness.  Write \(\tau=(1-t)a+tb\), where \(0<t<1\).
Concavity gives
\(u_\tau\geq(1-t)V_\theta+t u_b\).
Positivity and multilinearity of nonpluripolar products, together with
mass monotonicity under decreasing singularities, give
\[
       \int_X\theta_{u_\tau}^{\,n}
                    \geq(1-t)^n\int_X\theta_{V_\theta}^{\,n}>0.
\]
This also explains the positive-mass observation immediately before
DXZ Lemma~4.5 and completes the proof.
\end{proof}

\begin{corollary}\label{cor:pseff-terminal-filtration-slice}
If the DXZ terminal slice
\(u_c=\lim_{\tau\uparrow c}u_\tau\) is not identically \(-\infty\), then
\[
                 u_c=P_\theta[u_c]=P_\theta[u_c]_{\mathcal I}.
\]
This holds whether its top nonpluripolar mass is positive or zero.
\end{corollary}

\begin{proof}
DXZ defines \(u_c=\lim_{\tau\uparrow c}u_\tau\).
If this limit is not identically \(-\infty\), it is a decreasing limit of
psh functions and hence is psh.
For every \(\tau<c\), monotonicity gives
\[
 P_\theta[u_c]\leq P_\theta[u_\tau]=u_\tau,\qquad
 P_\theta[u_c]_{\mathcal I}\leq
                          P_\theta[u_\tau]_{\mathcal I}=u_\tau.
\]
Letting \(\tau\uparrow c\), and using the reverse inequalities from
\(u_c\leq0\), shows that this particular endpoint is also fixed by both
envelopes, even if its top mass is zero.
If \(u_c\equiv-\infty\), it is simply the extended-real terminal slice.
For \(\tau>c\) the slice is defined to be \(-\infty\).
These observations contain no reconstruction statement for unspecified
zero-mass filtration data.
\end{proof}

The corollary uses only the particular left-limit convention of DXZ.  It is
not a reconstruction theorem for an arbitrary zero-mass potential and does
not decide the separate zero-mass reconstruction question retained after
\Cref{prob:frontier-mp-pseff}; in the atomic inventory that question is
P14.8.

\section{Big non-nef classes and the zero-volume boundary}
\label{sec:pseff-boundary}

The two boundary phenomena are different.  A big non-nef class has
positive class volume but fails the generic-fibre nefness test.  A
zero-volume class has no normalized top-degree mass at all.

\begin{proposition}\label{prop:pseff-big-nonnef-surviving-core}
Assume that $\alpha$ is big and non-nef, and put
$V:=\vol(\alpha)>0$.  Then the unchanged MP relatively nef model cone
with generic-fibre class $\alpha$ is empty.  Nevertheless, every
$\boldsymbol\Gamma\in\PSH^{\NA}(X,\theta)$ has the canonical injective
ideal-valuative realization of
\Cref{thm:pseff-perturbation-independent-transform,thm:pseff-divisorial-inversion}.

If $\vol(\boldsymbol\Gamma)>0$, then
$\boldsymbol\Gamma$ is represented by a unique unperturbed
positive-mass $\mathcal I$-model test curve $\Gamma$. Put
$m=\vol(\Gamma)=M_\theta(\Gamma_{-\infty})>0$. Its normalized
\emph{relative} radial energy is
\[
 E_\theta^{\mathrm{NA,rel}}(\Gamma)
 :=c+\frac1{\vol(\Gamma)}\int_{-\infty}^{c}
   \left(
    \int_X\left\langle
       (\theta+\ddc\Gamma_\tau)^n
    \right\rangle-\vol(\Gamma)
   \right)\rd\tau.
\tag{71.49}\label{eq:pseff-big-dxz-energy}
\]
The constant subtracted inside this relative formula is $m$, the mass
of the \emph{initial} level $\Gamma_{-\infty}$, rather than the class
volume $V$. It is the barycentre of the probability measure
$m^{-1}\DHm(\Gamma)$, so
$E_\theta^{\mathrm{NA,rel}}(\Gamma+C)
  =E_\theta^{\mathrm{NA,rel}}(\Gamma)+C$.
In the book's unnormalized convention it equals
$m^{-1}\mathbf E^{\Gamma_{-\infty}}(\Gamma)$.

The \emph{class-normalized} DXZ energy in
\cite[formula~(43)]{DXZ23} instead subtracts and divides by $V$.
When $m=V$, it equals the displayed relative functional. When $m<V$,
its integrand tends to $m-V<0$ as $\tau\to-\infty$, so that source
functional equals $-\infty$, even if the relative energy is finite.
This distinction is essential for the finite-energy correspondence above.
The source defines its class-normalized energy to be $-\infty$ outside
the embedded unperturbed positive-mass locus.  None of these facts constructs
an MP psh cone, an MP Monge--Amp\`ere measure, or an MP energy in the
non-nef class.
\end{proposition}

\begin{proof}
The emptiness statement is
\Cref{prop:pseff-relative-nef-obstruction}: any member of the unchanged
model cone would restrict to a relatively nef class and force $\alpha$
to be nef.  The transform and its injectivity are independent of nefness
by
\Cref{thm:pseff-perturbation-independent-transform,thm:pseff-divisorial-inversion}.

The positive-volume recovery and uniqueness are
\cite[Lemma~3.10]{Xia23Operations}, restated in
\Cref{thm:tna-positive-volume,cor:tna-positive-no-ambiguity}.  For a big
class, the relative energy definition \eqref{eq:rays-test-energy},
divided by the positive initial mass $m$, is exactly
\eqref{eq:pseff-big-dxz-energy}. The intrinsic DH moment formula
\Cref{thm:ok-dh-flag-independent} identifies this value with the normalized
barycentre, including the value $-\infty$. For bounded curves this is
also \Cref{thm:cma-test-curve-moments}. The class-normalized formula and
its $-\infty$ extension are \cite[formula~(43)]{DXZ23}.
The interior masses tend to the initial mass as $\tau\to-\infty$,
so replacing $m$ by a strictly larger $V$ gives the asserted
divergence. Translation leaves the initial mass unchanged and translates
the DH variable, proving the relative translation formula.  The last sentence follows because no
step changes the generic-fibre class in the MP model signature or defines
intersection measures for a non-nef cohomological model.
\end{proof}

\begin{proposition}\label{prop:pseff-zero-volume-boundary}
Assume that $\vol(\alpha)=0$.  Then every
$\varphi\in\PSH(X,\theta)$ satisfies
\[
 \int_X\left\langle
       (\theta+\ddc\varphi)^n
 \right\rangle=0.
\tag{71.50}\label{eq:pseff-zero-top-mass}
\]
Consequently, neither the probability normalization
$V^{-1}\langle(\theta+\ddc\varphi)^n\rangle$ nor the normalized energy
formula \eqref{eq:pseff-big-dxz-energy} is defined when $V=0$.
The inverse-system space and the injective transform
\eqref{eq:pseff-canonical-transform-map} still make sense.  The statement
does not rule out a future unnormalized or lower-dimensional boundary
theory.
\end{proposition}

\begin{proof}
Let $V_\theta$ have minimal singularities.  It is less singular than
$\varphi$.  The unrestricted monotonicity theorem for total
nonpluripolar mass \cite[Theorem~1.1]{DDNL18mono} gives
\[
 0\leq\int_X\left\langle
       (\theta+\ddc\varphi)^n
      \right\rangle
 \leq\int_X\left\langle
       (\theta+\ddc V_\theta)^n
      \right\rangle
 =\vol(\alpha)=0,
\]
which proves \eqref{eq:pseff-zero-top-mass}.  Division by $V=0$ is
undefined, proving the normalization statements.  DXZ define
$\PSH^{\NA}(X,\theta)$ as the inverse limit over the big classes
$\theta+\omega$, so its definition does not divide by $V$.
The construction and proof of
\Cref{thm:pseff-perturbation-independent-transform} likewise use only
multiplier ideals and endpoints.  This proves the last assertion.
\end{proof}

\begin{example}[A nef class of numerical dimension one]
\label{ex:pseff-nef-zero-volume}
Let $X=C_1\times C_2$ be a product of compact Riemann surfaces, and let
$\omega_1$ be a K\"ahler form on $C_1$.  The class
$\alpha=\{\operatorname{pr}_1^*\omega_1\}$ is nef and nonzero, but
\[
                 \alpha^2=0,
                 \qquad \vol(\alpha)=0.
\tag{71.51}\label{eq:pseff-product-zero-volume}
\]
Thus the obstruction in
\Cref{prop:pseff-zero-volume-boundary} is not confined to the zero class.
The inverse-system transform still records divisorial singularity data,
but a probability Monge--Amp\`ere measure cannot be obtained by dividing
its top mass by $\vol(\alpha)$.
\end{example}

\begin{example}[An effective exceptional boundary class]
\label{ex:pseff-exceptional-zero-volume}
On the blowup of $\mathbf P^2$ at a point, the exceptional class
$\alpha=\{[E]\}$ is pseudoeffective and satisfies
$\vol(\alpha)=0$.  It is not nef because $E^2=-1$.  It therefore lies
on both sides of the boundary ledger: the unchanged relatively nef MP
model signature is excluded by
\Cref{prop:pseff-relative-nef-obstruction}, while the normalized
top-degree formulas are excluded by
\Cref{prop:pseff-zero-volume-boundary}.  The DXZ inverse system retains
the multiplier ideals of the divisorial pole along $E$, and
\Cref{thm:pseff-divisorial-inversion} recovers that pole profile from the
values on the ray $\mathbf Q_{>0}\ord_E$.
\end{example}

\section{Exact status and remaining boundaries}
\label{sec:pseff-status}

The conclusions of this chapter can be read as the following ledger.
\begin{center}
\small
\begin{tabular}{p{0.29\textwidth}p{0.62\textwidth}}
\toprule
Question & Audited status \\
\midrule
All pseudoeffective classes
& The DXZ inverse system has the canonical ideal-valuative transform
  \eqref{eq:pseff-canonical-transform-map}.  Its divisorial restriction is
  the sourced DXZ analytification, which is injective and order-reflecting. \\
Big non-nef classes
& The transform and normalized relative radial energy survive; the
  class-normalized DXZ energy is finite only on its full-initial-mass
  finite-energy locus. The unchanged
  relatively nef MP model signature is empty, so no MP measure or energy
  is inferred. \\
Zero class volume
& The transform survives, while all top nonpluripolar masses vanish and
  probability normalization is unavailable. \\
K\"ahler classes
& The raw transform is an order isomorphism from the full DXZ cone to the
  full MP cone by \Cref{thm:p143-full-kahler-comparison}.  On the
  finite-energy locus, MP and DXZ objects are in bijection through their common maximal
  radial asymptotic data.  An initial potential selects the unique
  maximal representative with that initial value; their energies agree
  after the normalization fixed above.  Their raw ideal-valuative
  transform agrees with the MP potential at every semivaluation by
  \Cref{cor:p143-all-point-kahler-comparison}, and its image is exactly
  $\mathcal E^1_{\mathrm{MP}}(\alpha)$ by
  \Cref{cor:p141-finite-energy-raw-image}. \\
Projective pseudoeffective real N\'eron--Severi classes
& The full DXZ and Boucksom--Jonsson spaces agree by the source's
  directed-limit identification.  For a $\mathbf Q$-line-bundle class
  this is the literal ideal-valuative formula at every semivaluation.  If
  the real numerical class is ample, the full MP and Boucksom--Jonsson
  cones agree by \Cref{thm:pseff-projective-full-mp-bj}; for an ample
  $\mathbf Q$-class, the full DXZ and MP cones therefore agree without
  an energy hypothesis. \\
Operations
& Translation holds on every semivaluation.  The displayed algebraic and
  monotone formulas are asserted pointwise only on rational divisorial
  valuations, with their endpoint hypotheses. \\
\bottomrule
\end{tabular}
\end{center}

The following boundaries remain open in the general compact K\"ahler
setting.
\begin{enumerate}
\item Give an intrinsic necessary-and-sufficient characterization of the
      image of $\mathfrak T_\theta$ inside the extended-real functions
      on $X^{\bethan}$, beyond the full K\"ahler characterization in
      \Cref{thm:p143-full-kahler-comparison} and the established projective
      regimes.
\item Construct a positivity theory for big non-nef classes that retains
      movable or $P$-equivalence data instead of silently replacing it
      by $\mathcal I$-equivalence.
\item Construct a meaningful unnormalized or lower-dimensional energy and
      measure theory at zero volume.
\end{enumerate}
The upper-semicontinuity statement formerly listed among these mathematical
boundaries is now \Cref{thm:p142-raw-transform-usc,cor:p142-pseudoeffective-raw-usc}.
These are genuine open boundaries.  No item in this list is used as an
input elsewhere in the chapter.


\paragraph{Source versions and energy conventions.}
For DXZ, the comparison and filtration locators refer to version~3 of
arXiv 2302.02541 (17 July 2024). Its class-normalized energy formula (43)
is formula (34) in the published 2025 paper; both subtract and divide by
the class volume. The relative functional in
\eqref{eq:pseff-big-dxz-energy} follows the book's relative radial energy
and DH first-moment formulas, rather than changing that source definition.
The operations locators follow Xia's arXiv 2312.17150, version~3 (22 October 2024),
and the MP locators follow arXiv 2409.06221, version~2 (18 September 2025). These version choices
preserve the stated source interfaces and the distinct energy domains.
